\chapter{실험 부록}
\label{sec:supplement}

여기에는 각 장마다 핵심 주장과 그 근거를 모았다. 그것을 측정한 실험, 정확히 무엇을 측정했는지, 어디에 발표되었는지이다. 오른쪽 열의 상태는 주장이 얼마나 단단한지를 말해 준다. \emph{측정됨}은 직접 실험이 있다는 뜻이다. \emph{이론}은 본문이나 고전적 논문에서 유도한 수학적 귀결이다. \emph{모델}은 이 책의 모델로 계산한 결과이다(스크립트는 \texttt{models/} 폴더에 있다). \emph{추정}은 직접 측정 없는 크기 정도의 수치이다. \emph{가설}은 그럴듯하지만 아직 실험으로 증명되지 않은 설명이다. 실험이 없는 곳에는 그렇다고 적었다.

\section{\ref{sec:neurontypes}장. 뉴런의 유형}

이 장의 수치는 사람의 뇌에서 세포를 직접 센 자료가 있는 곳에서는 그것에 기댄다. “뉴런 하나, 신경전달물질 하나” 규칙은 세포 아틀라스와 예외까지 찾아낸 실험에, \nt{DOPA}의 역할은 스파이크 기록에, 나머지 브로드캐스트 채널의 역할은 가설에 기댄다.

{\footnotesize
\setlength{\tabcolsep}{3pt}\renewcommand{\arraystretch}{1.15}
\begin{longtable}{>{\raggedright}p{0.5cm}>{\raggedright}p{4.0cm}>{\raggedright}p{5.6cm}>{\raggedright}p{2.3cm}>{\raggedright\arraybackslash}p{1.7cm}}
\toprule
No. & 명제 & 실험과 측정한 것 & 출처 & 상태 \\
\midrule\endfirsthead
\toprule
No. & 명제 & 실험과 측정한 것 & 출처 & 상태 \\
\midrule\endhead
1 & 사람의 뇌에는 뉴런이 약 $8.6\cdot10^{10}$개 있고, \nt{GLUT} 뉴런의 거의 전부는 소뇌의 작은 뉴런이다(표~\ref{tab:types}) & 등방성 분획법: 조직을 녹여 세포핵의 균일한 현탁액으로 만들고, 뉴런 표지 NeuN을 가진 핵을 센다. 성인 남성의 뉴런은 $86.1\pm8.1$\,십억(billion) 개이다. 그중 대뇌 피질에 있는 것은 $19\,\%$뿐이고 대부분은 소뇌에 있다 & \cite{Azevedo2009} & 측정됨 \\
2 & 거의 모든 뉴런에는 주된 신경전달물질이 하나 있지만(명제~\ref{prop:dale}) 예외도 있다 & 생쥐 뇌의 전사체 아틀라스: 세포 약 $4\cdot10^6$개, 클러스터 $5322$개. 합성과 포장 유전자에 따라 클러스터에 신경전달물질을 배정했다. 예외는 직접 측정되었다. 절편에서 \nt{DOPA} 뉴런은 같은 소포 수송체를 통해 GABA도 내보내 선조체 뉴런을 억제한다. 일부 \nt{DOPA} 축삭에서는 글루탐산과 도파민이 같은 전선의 서로 다른 말단에서 나온다(전자현미경과 광유전학) & \cite{Strata1999,Yao2023,Tritsch2012,Zhang2015} & 측정됨 \\
3 & 가중치 행렬의 열 전체가 같은 부호이다~\eqref{eq:dalesign} & 명제~\ref{prop:dale}와, 글루탐산 수용체는 거의 모두 흥분성이고 GABA 수용체는 억제성이라는 사실에서 장에서 유도한 것이다. “한 뉴런의 모든 수신자가 같은 부호의 가중치를 받는다”를 일반 규칙으로 직접 검증한 연구는 없다. 반례는 \nt{DOPA} 뉴런의 GABA 공동 방출(2행)과 부호가 다른 수용체를 가진 수신자들(11행)이다 & 이 장의 유도 & 이론 \\
4 & 피질 뉴런의 4분의 3에서 5분의 4는 \nt{GLUT}, 5분의 1에서 4분의 1은 \nt{GABA}이다 & 원숭이 피질 $10$개 영역에서 피질 전체 두께에 걸친 폭 $50$\,\textmu m의 기둥을 GABA로 면역염색했다. $8$개 영역에서 GABA 뉴런은 전체 뉴런의 약 $25\,\%$, 일차 시각 피질에서는 $20\,\%$ 미만이다 & \cite{Hendry1987} & 측정됨 \\
5 & \nt{GLUT} 뉴런의 출력은 약 $10^4$개이다 & 부검 입체해석학: 젊은 남성의 신피질에 시냅스 $1.64\cdot10^{14}$개(전자현미경, 디섹터법)와 뉴런 약 $2.3\cdot10^{10}$개. 비를 구하면 뉴런당 시냅스 $\sim7\cdot10^3$개이다. 평균적으로 입력 수는 출력 수와 같다 & \cite{Tang2001synapses,Pakkenberg1997} & 측정됨 \\
6 & \nt{DOPA} 뉴런의 출력은 말단 $10^5$--$10^6$개로 가지를 친다. 이것은 표적을 덮는 범위에서 얻은 추정치이지 센 값이 아니다 & 쥐의 \nt{DOPA} 뉴런 하나하나를 바이러스로 막 표지하고 축삭을 완전히 재구성했다. 뉴런 하나가 선조체 부피의 $0.45$에서 $5.7\,\%$(평균 $2.7\,\%$)를 덮는다. 측정한 것은 덮는 범위이지 말단 수가 아니다. 말단 수는 덮는 범위에서 자릿수 수준으로 유도했고, 말단을 직접 센 자료는 없다. \nt{SERO}와 \nt{NORA}의 말단 수는 거친 추정치이다 & \cite{Matsuda2009} & 측정됨 \\
7 & 브로드캐스트 뉴런은 적다. \nt{DOPA} $\sim5\cdot10^5$(두 무리 중 큰 쪽, 하한), \nt{SERO} $\sim3\cdot10^5$(두 부분 집계의 합), \nt{HIST} $\sim6\cdot10^4$(표~\ref{tab:types}, 그림~\ref{fig:typepie}) & 사람에서의 입체해석학적 집계: 흑색질의 색소 뉴런 $550\,000$개(배쪽 피개 영역 제외); 등쪽 솔기핵의 뉴런 $235\,000$개(핵의 모든 뉴런)와 다리뇌 피개의 세로토닌 뉴런 $125\,000$개; 시상하부의 히스타민 뉴런 약 $64\,000$개. \nt{NORA}, \nt{GLYC}, \nt{ACET}, \nt{ADRE}는 직접 센 자료가 없어 표에는 거친 자릿수 추정치를 넣었다 & \cite{Pakkenberg1991,Baker1990,Baker1991,Airaksinen1991} & 측정됨 \\
8 & 틈 속 글루탐산은 약 1밀리몰까지 치솟았다가 $\sim1\ms$ 만에 줄어든다 & 시냅스 전달 중 빠르게 떨어지는 경쟁적 차단제가 NMDA 수용체에서 밀려나는 정도로 측정했다(해마 뉴런 배양). 정점 $1.1$\,mM, 감소 시간 상수 $1.2\ms$ & \cite{Clements1992} & 측정됨 \\
9 & 작동 중인 피질에서는 흥분 전류와 억제 전류가 거의 균형을 이루며, 뉴런은 문턱값 근처에 머문다 & 이론: 연결이 강한 네트워크는 스스로 균형 체제에 이른다. 실험: 흰담비 전전두 피질 in vivo에서 “업” 상태 동안 시냅스 전류의 역전 전위는 수백 밀리초 동안 약 $-37$\,mV로 유지되는데, 컨덕턴스는 $\sim21$\,nS만큼 변한다. 흥분과 억제가 함께 오르내린다 & \cite{vanVreeswijk1996,Haider2006} & 측정됨 \\
10 & \nt{DOPA} 뉴런은 보상 예측 오차~\eqref{eq:rpe}를 알린다(그림~\ref{fig:rpe}) & 원숭이 \nt{DOPA} 뉴런의 스파이크: 예고 없는 주스에 버스트; 학습 뒤에는 신호에 버스트, 예측된 주스에는 무응답; 약속한 주스를 주지 않으면 발화율이 움푹 꺼진다. 정량 실험에서 발화율은 현재 보상과 과거 보상의 가중 평균의 차에 비례하지만, 양의 오차에서만 그렇다. 식~\eqref{eq:rpe} 자체는 시간차 알고리즘이다 & \cite{Schultz1997,Bayer2005,SuttonBarto2018} & 측정됨 \\
11 & 부호와 속도는 수용체가 정한다. 이온성은 1밀리초 미만, 대사성은 훨씬 느리다(명제~\ref{prop:receptor}) & 해마 뉴런의 막 조각에 글루탐산을 짧게 가했다. 이온성 수용체를 통한 전류는 $0.2$--$0.6\ms$ 만에 ($20$에서 $80\,\%$까지) 오르고 $2$--$3\ms$ 만에 줄어든다. 피라미드 뉴런의 느린 억제 전위는 대사성 GABA 수용체의 선택적 차단제로 사라진다. 도파민 수용체에는 작용이 반대인 두 계열이 있다. 선조체에서 D1 수용체를 가진 뉴런은 시냅스 강화에, D2를 가진 뉴런은 시냅스 약화에 도파민이 필요하다 & \cite{Colquhoun1992,DutarNicoll1988,Beaulieu2011,Shen2008} & 측정됨 \\
12 & 브로드캐스트 채널은 전선 한 가닥이 아니라 적은 수의 선로로 된 다발이다(\ref{sec:buses}절) & 생쥐 등쪽 솔기핵 세로토닌 뉴런의 바이러스-유전적 표지: 편도체로 출력을 보내는 뉴런과 이마엽 피질로 보내는 뉴런은 핵 안의 다른 자리에 있고, 서로 보완하는 영역으로 가지를 치며, 불쾌한 자극에 반대로 응답한다. 총설: 청반(\nt{NORA}) 뉴런의 하위 무리들은 서로 다른 과정을 부호화한다 & \cite{Ren2018,Poe2020} & 측정됨 \\
13 & \nt{NORA}는 이득 $g$를 정한다 & 적응적 이득 가설; 카테콜아민이 전달 특성의 기울기를 높이는 네트워크 모델. “노르아드레날린과 함께 $g$가 커진다”를 네트워크 전체에서 직접 측정한 자료는 없다. 원숭이에서 동공 지름이 청반 뉴런의 스파이크를 따라간다는 것은 측정되었다 & \cite{AstonJones2005,ServanSchreiber1990,Joshi2016} & 가설 \\
14 & \nt{ACET}은 “쓰기/읽기”를 전환하고 학습률을 정한다 & 가설. 근거: 쥐 후각 피질 절편에서 아세틸콜린 작용제($100$\,\textmu M)는 내부의 “회상” 시냅스를 $60\,\%$ 넘게 약화하지만 입력 시냅스는 $15\,\%$ 미만으로 약화한다 & \cite{Hasselmo2006,HasselmoBower1992} & 가설 \\
15 & \nt{SERO}는 할인율 $\gamma$를 정한다; 세로토닌 뉴런은 초당 스파이크 몇 개로 고르게 작동하고 수면의 한 단계에서는 거의 잠잠하다 & “세로토닌 $=\gamma$” 가설. 가장 강력한 근거: 생쥐 등쪽 솔기핵 세로토닌 뉴런을 광유전학으로 활성화하면 지연된 보상을 기다리다 일찍 포기하는 횟수가 줄어든다. 발화율과 급속 안구 운동 수면 중의 침묵은 측정되었다(기록 총설) & \cite{Doya2002,Miyazaki2014,JacobsAzmitia1992} & 가설 \\
\bottomrule
\end{longtable}}

\section{\ref{sec:glutcomputer}장. \nt{GLUT} 뉴런은 무엇을 계산하는가}

이 장의 논리적 명제는 음이 아닌 가중치로 된 회로에 관한 정리로, 장 안에서 유도했다. 실험은 뉴런 모델을 뒷받침하고, 이 정리들에 따라 조립한 회로(동시성 검출기, 지연선, 자기 유지 무리, 사슬)가 뇌에 실제로 있다는 것을 보여 준다.

{\footnotesize
\setlength{\tabcolsep}{3pt}\renewcommand{\arraystretch}{1.15}
\begin{longtable}{>{\raggedright}p{0.5cm}>{\raggedright}p{4.0cm}>{\raggedright}p{5.6cm}>{\raggedright}p{2.3cm}>{\raggedright\arraybackslash}p{1.7cm}}
\toprule
No. & 명제 & 실험과 측정한 것 & 출처 & 상태 \\
\midrule\endfirsthead
\toprule
No. & 명제 & 실험과 측정한 것 & 출처 & 상태 \\
\midrule\endhead
1 & 뉴런은 문턱값과 초기화가 있는 RC 회로이다~\eqref{eq:glutneuron} & 쥐 피질 피라미드 뉴런(절편)의 세포체에 살아 있는 뇌의 시냅스 흐름과 비슷한 잡음 전류를 주입했다. 평균 발화율이 전류의 평균과 산포에 어떻게 의존하는지는 느린 발화율 적응을 더한 누설 적분 발화 뉴런으로 기술된다. $\tau$와 지연의 범위는 장에서 실험 인용 없이 제시했다 & \cite{Rauch2003} & 측정됨 \\
2 & 모델~\eqref{eq:glutneuron}은 단순화이다. 입력 가지가 스스로 국소 스파이크를 낸다 & 생쥐 시각 피질 피라미드 뉴런의 가지돌기에서 in vivo로 직접 기록: 시각 자극이 NMDA 수용체에 의존하는 국소 가지돌기 스파이크를 일으키며, 이를 억누르면 뉴런의 방위 조율이 넓어진다 & \cite{Smith2013} & 측정됨 \\
3 & 일치 창 $\Delta_{\max}=\tau\ln\frac{w}{\theta-w}$~\eqref{eq:window}; $\theta=1.5w$이면 $0.7\tau$ & 모델~\eqref{eq:glutneuron}의 결과이다. 좁은 합산 창을 직접 측정한 자료는 빠른 억제가 있는 뉴런에 대해 있다(\ref{sec:gaba}장, 그 장 표의 3행) & 이 장의 유도 & 이론 \\
4 & \nt{GLUT} 뉴런만으로 된 네트워크는 입력의 단조 함수만 계산한다(명제~\ref{prop:monotone}) & 장에서의 증명: 우변이 감소하지 않는 침묵에서의 반복. 이것은 모델에 관한 명제이다. 억제 없는 살아 있는 네트워크에서 이를 검증한 실험은 없다 & 이 장의 유도 & 이론 \\
5 & 감각 기관은 전선 쌍을 준다. 어떤 세포는 자극이 강해질 때, 다른 세포는 약해질 때 응답한다 & 망막: 이미 쌍극 세포에서 원뿔 세포의 신호가 켜짐 채널과 꺼짐 채널로 갈라진다. 원숭이에서 켜짐 쌍극 세포를 약리학적으로 차단: 채널은 피질까지 따로 가며, 차단 뒤에는 빛이 강해지는 것은 잘 검출하지 못하지만 약해지는 것은 그렇지 않다 & \cite{Schiller1992} & 측정됨 \\
6 & 이중 레일 부호가 있으면 \nt{GLUT} 뉴런 네트워크는 공통 박자 없이 어떤 논리 함수든 비동기적으로 계산하며, 대가는 두 배의 뉴런 수이다(명제~\ref{prop:dualrail}, \eqref{eq:dualrail}, 그림~\ref{fig:dualxor}) & 이중 레일 부호와 신호 자체로 준비 상태를 판정하는 것은 지연에 둔감한 비동기 회로의 표준 기법이다. 뉴런 여섯 개짜리 배타적 OR 회로는 장에서 조립했다. 뇌가 논리 함수를 바로 이중 레일 부호로 계산한다는 것을 보여 주는 실험은 없다 & \cite{Sparso2001} & 이론 \\
7 & 사람에서 두 귀에 소리가 도착하는 시간의 최대 차는 $\approx0.6\ms$이고, 뇌는 약 10마이크로초를 구별한다 & 두 선택지 실험의 청취자: 시간 차 변별 문턱(정답 75\,\%)은 $150$--$1700$\,Hz 대역 잡음에서 $9$\,\textmu s, $1$\,kHz 순음에서 $11$\,\textmu s, 딸깍 소리에서 $28$\,\textmu s. 상한 $0.6\ms$는 장에서 $0.22\,\text{m}/343\,\text{m/s}$로 계산했다 & \cite{KlumppEady1956} & 측정됨 \\
8 & 새에서는 지연선과 동시성 검출기가 시간 차를 위치로 바꾼다~\eqref{eq:itd}(그림~\ref{fig:itd}) & 원숭이올빼미: 두 귀에서 온 축삭이 동시성 검출기 핵에 반대쪽에서 들어온다. 축삭을 따라 스파이크 도착 시간이 깊이에 따라 규칙적으로 변하고, 핵을 가로지르는 지연의 폭($700$\,\textmu m에 걸쳐 $160$\,\textmu s)은 두 귀 사이 지연의 범위와 일치한다 & \cite{Carr1990} & 측정됨 \\
9 & 포유류에서는 지연의 열이 발견되지 않았다. 같은 문제를 \nt{GLYC} 뉴런의 정확한 시간의 억제가 있는 동시성 검출기가 푼다 & 게르빌루스쥐 내측 상올리브에서 in vivo 기록: 응답이 방향 지도를 이루지 않는다. 미세 피펫으로 글리신과 그 차단제를 가하면, 시간을 정확히 맞춘 \emph{글리신} 억제가 지연의 작동 범위를 정한다는 것이 드러난다. 여기서 억제는 \nt{GABA}가 아니라 \nt{GLYC}에서 온다 & \cite{Brand2002,Grothe2010} & 측정됨 \\
10 & 되먹임이 강한 무리는 플립플롭이다. 자기 유지 활동은 자극 뒤 몇 초 동안 지속된다(그림~\ref{fig:bistable}) & 기억한 점으로 눈을 늦게 움직이는 과제(지연 $1$--$6$\,s)를 하는 원숭이의 전전두 피질: 과제와 관련된 뉴런 $170$개 중 $87$개가 지연 동안 발화율을 바꾸며, 그중 $79\,\%$에서 이 변화는 기억한 점의 방향에 의존한다. 이 활동을 두 안정 상태를 가진 되먹임이 유지한다는 것은 이론이다 & \cite{Funahashi1989,Wang2001} & 측정됨 \\
11 & 유입이 $\approx0.7\tau$보다 길면 비트가 써진다(그림~\ref{fig:recurrentgroup}b) & 방정식 $\tau\,\dd r/\dd t=-r+f(wr+I)$를 $w=1$, $I=0.6$에서 오일러 방법으로 계산했다. \texttt{models/}에 스크립트는 없다. 실험은 없다 & 이 장의 계산 & 모델 \\
12 & 기억은 스파이크 없이 거의 시냅스 촉진만으로 저장할 수 있다 & 이론: 말단의 잔류 칼슘이 기록을 약 1초 유지한다. 근거: 흰담비 내측 전전두 피질(여러 뉴런 동시 기록)에 여운이 긴 촉진성 시냅스로 연결된 피라미드 뉴런 하위 네트워크가 있다. 기억 유지 중 “조용한” 구간 관찰의 총설. 피질이 언제 어떤 방법을 쓰는지는 열린 문제이다 & \cite{Mongillo2008,WangMarkram2006,Stokes2015} & 가설 \\
13 & 기억은 피로로 지워진다. 신경전달물질 저장량이 고갈된다 & 피질 피라미드 뉴런의 쌍 기록: 스파이크가 잦으면 시냅스 응답이 줄고, 감소 속도는 방출 확률로 정해진다. 바로 이것이 자기 유지 무리를 끈다는 것은 직접 보이지 않았다 & \cite{Tsodyks1997} & 측정됨 \\
14 & 무리의 사슬은 사건이 시동하는 타이머이다. 명금류에서는 뉴런들이 엄격히 차례대로 발화한다 & 잠잘 때와 노래할 때 금화조의 고위 발성 중추에서 식별된 뉴런을 기록: 운동핵으로 출력을 보내는 뉴런은 저마다 노래의 정확한 한 순간에 짧은 버스트 하나를 내고, 뉴런들은 차례로 발화한다 & \cite{Hahnloser2002} & 측정됨 \\
\bottomrule
\end{longtable}}

\section{\ref{sec:gaba}장. \nt{GLUT}와 \nt{GABA}의 협업}

이 장의 논리적, 동역학적 명제는 선형 해석과 옴의 법칙으로 장 안에서 유도했다. 실험은 이 명제들이 기대는 회로, 즉 지연된 금지, 흥분 뒤에 따라오는 직접 억제, 억제에 의한 안정화, \nt{GLUT}--\nt{GABA} 루프의 리듬이 뇌에서 실제로 작동한다는 것을 보여 준다.

{\footnotesize
\setlength{\tabcolsep}{3pt}\renewcommand{\arraystretch}{1.15}
\begin{longtable}{>{\raggedright}p{0.5cm}>{\raggedright}p{4.0cm}>{\raggedright}p{5.6cm}>{\raggedright}p{2.3cm}>{\raggedright\arraybackslash}p{1.7cm}}
\toprule
No. & 명제 & 실험과 측정한 것 & 출처 & 상태 \\
\midrule\endfirsthead
\toprule
No. & 명제 & 실험과 측정한 것 & 출처 & 상태 \\
\midrule\endhead
1 & 피질 뉴런의 5분의 1에서 4분의 1이 \nt{GABA} 뉴런이다 & 원숭이 피질 $10$개 영역의 GABA 면역염색: $8$개 영역에서 뉴런의 약 $25\,\%$, 일차 시각 피질에서는 $20\,\%$ 미만. 피질 억제성 뉴런의 다양성에 관한 총설 & \cite{Hendry1987,Markram2004} & 측정됨 \\
2 & 억제가 하는 일은 붙는 자리가 크게 좌우한다. 가지돌기, 세포체, 스파이크가 생기는 자리, 다른 뉴런 전선의 말단 & 총설: 피질 \nt{GABA} 뉴런의 부류는 수신자 위의 표적 자리로 구별된다. 피라미드 뉴런의 세포체와 가지돌기에서 동시 기록: 직접 억제는 가지돌기보다 세포체에서 훨씬 강하다. 척수에서 다른 뉴런 전선의 말단에 가해지는 억제는 입력 선로 하나의 신경전달물질 방출을 줄인다(측정 총설) & \cite{Markram2004,Pouille2001,Rudomin1999} & 측정됨 \\
3 & 중개자를 거친 부호 전환의 값은 지연 하나, 약 1밀리초이다. 흥분 뒤에 따라오는 억제는 일치 창을 좁힌다 & 쥐 해마 절편의 피라미드 뉴런: 중개자 하나를 거친 억제가 직접 흥분 뒤에 짧은 지연으로 따라오며, 두 흥분성 입력이 문턱값까지 합산되는 창은 $2\ms$보다 짧다. 이 억제가 약한 가지돌기에서는 창이 넓다 & \cite{Pouille2001} & 측정됨 \\
4 & 금지 $a\wedge\bar b$~\eqref{eq:veto}는 OR, 상수 1과 함께 함수적으로 완전하다(명제~\ref{prop:gabacomplete}) & 장에서의 증명. 고전적 결과: 억제성 입력이 있는 문턱 소자의 네트워크는 어떤 논리식이든 구현한다. 모델에 관한 명제이며, 실험은 없다 & \cite{McCullochPitts1943} & 이론 \\
5 & 지연된 금지는 최적 속도 $v^*=\Delta x/d$~\eqref{eq:vstar}의 운동 방향 검출기를 준다 & 토끼 망막: 방향 선택성은 “무효” 방향으로 움직일 때 흥분을 끄는 억제로 만들어진다. 선택성 세포의 흥분성 입력과 억제성 입력을 따로 기록한 결과, 별 모양 아마크린 세포(\nt{GABA})가 “무효” 쪽을 향한 돌기로만 비대칭적으로 억제한다는 것이 드러났다. $v^*$ 공식은 장에서 유도했다 & \cite{Barlow1965,Fried2002} & 측정됨 \\
6 & 분로 억제는 전압에서 빼지 않고 전압을 나눈다~\eqref{eq:shunt}(그림~\ref{fig:shunt}) & 염소 평형 전위가 휴지 전위 근처일 때 세 컨덕턴스로 된 분배기에 대한 옴의 법칙. 스파이크를 내지 않는 뉴런에 관한 것이다 & 이 장의 유도 & 이론 \\
7 & 분로는 스파이크 발화율에 뺄셈으로 작용하고, 나눗셈은 억제와 균형 잡힌 배경 잡음의 조합에서 나온다 & 모델: 발화하는 뉴런에서는 분로를 지나는 전류가 거의 일정하여 발화율에 대한 효과가 뺄셈이다. 실험: 쥐 체성감각 피질(절편)의 피라미드 뉴런에 동적 클램프로 흥분성, 억제성 배경 컨덕턴스를 주입했다. 둘을 동시에 균형 있게 늘리면 발화율의 뚜렷한 이동 없이 응답 기울기가 나누어진다. 총설: 총활동으로 나누는 연산은 뇌의 여러 부위에서 나타난다 & \cite{HoltKoch1997,Chance2002,Carandini2012} & 측정됨 \\
8 & $\beta>1$이면 상호 억제가 승자 하나를 고른다(명제~\ref{prop:wta}, \eqref{eq:wta}) & 고윳값 $-(1\pm\beta)/\tau$는 장에서 유도했다. $\beta=0.5$에서 발화율 $0.73$과 $0.53$은 고정점 $r_{1,2}=(I_{1,2}-\beta I_{2,1})/(1-\beta^2)$이다. \texttt{models/}에 스크립트는 없다. 장에서 직접 실험을 제시하지 않는다 & 이 장의 유도 & 이론 \\
9 & \nt{GABA}를 거친 신호로 기억을 초기화; 상호 억제하는 두 무리는 래치이다 & 그림~\ref{fig:bistable}과 명제~\ref{prop:wta}에서 나온다. 실험은 없다 & 이 장의 유도 & 이론 \\
10 & \nt{GLUT}--\nt{GABA} 루프의 평형은 억제가 충분히 강하고 충분히 빠르면 안정하다(명제~\ref{prop:ei}) & 두 집단 모델~\eqref{eq:ei}. 조건 (i)과 (ii)는 그 행렬의 행렬식과 대각합이며 장에서 유도했다 & \cite{WilsonCowan1972} & 이론 \\
11 & $w_{EE}=1.5$, $w_{EI}w_{IE}=2$에서 빠른 억제($\tau_I=2\ms$)는 $E^*=0.67h$를 유지하고, 느린 억제($30\ms$)는 진동을 부추긴다(그림~\ref{fig:ei}) & \eqref{eq:ei}의 수치해, 스크립트 \texttt{figures/make\_ei.py} & 계산 & 모델 \\
12 & 피질은 억제 안정화 체제에서 작동한다. 특징은 \nt{GABA} 뉴런을 밀어 주면 그 발화율이 떨어진다는 것이다~\eqref{eq:isn} & 이론이 역설적 응답을 예측했다. 실험: 고양이 일차 시각 피질의 세포내 기록에서 주변 자극으로 응답이 억눌릴 때 억제성 컨덕턴스는 늘지 않고 흥분성 컨덕턴스와 함께 줄어든다. 생쥐의 시각, 체성감각, 운동 피질에서 PV 뉴런을 광유전학으로 흥분시키면 자극이 없을 때와 가벼운 마취에서도 억제성 뉴런의 발화율이 낮아진다. 그림~\ref{fig:ei}의 수치에서 계수 $-1/3$은 장에서 유도했다 & \cite{Tsodyks1997isn,Ozeki2009,Sanzeni2020} & 측정됨 \\
13 & “\nt{GLUT}가 \nt{GABA}를 흥분시키고, \nt{GABA}가 \nt{GLUT}를 억제한다”는 루프는 주기 $12$--$50\ms$의 빠른 리듬을 낳는다 & 생쥐 체성감각 피질에서 빠른 \nt{GABA} 뉴런을 $8$--$200$\,Hz로 광유전학 자극하면 감마 리듬($20$--$80$\,Hz, 주기 $12$--$50\ms$)이 선택적으로 강해진다. \nt{GLUT} 뉴런 자극은 더 느린 리듬만 강화한다 & \cite{Cardin2009,Buzsaki2012} & 측정됨 \\
\bottomrule
\end{longtable}}

\section{\ref{sec:code}장. 모스 부호}

이 장의 한계 추정은 정보 이론과 계산이다. 채널의 어디서 잡음이 생기는지, 뇌가 얼마나 빨리 작동하는지, 스파이크 하나가 얼마인지는 측정되었다. 피질이 실제로 어떤 부호를 쓰는지는 열린 문제로 남아 있다.

{\footnotesize
\setlength{\tabcolsep}{3pt}\renewcommand{\arraystretch}{1.15}
\begin{longtable}{>{\raggedright}p{0.5cm}>{\raggedright}p{4.0cm}>{\raggedright}p{5.6cm}>{\raggedright}p{2.3cm}>{\raggedright\arraybackslash}p{1.7cm}}
\toprule
No. & 명제 & 실험과 측정한 것 & 출처 & 상태 \\
\midrule\endfirsthead
\toprule
No. & 명제 & 실험과 측정한 것 & 출처 & 상태 \\
\midrule\endhead
1 & 피질 \nt{GLUT} 뉴런의 발화율은 초당 1개 미만에서 수십 개이며, 뉴런은 대부분의 시간을 아래쪽 경계 근처에서 작동한다 & 깨어 있는 쥐의 청각 피질에서 세포에 붙인 피펫으로 기록(뉴런 선택이 활동과 무관): 어느 순간이든 초당 $20$개가 넘는 발화율을 내는 뉴런은 $5\,\%$ 미만이다. 일부 뉴런의 배경 발화율은 초당 $0.01$보다 낮다. 발화율 분포는 로그 정규 분포이다 & \cite{Hromadka2008} & 측정됨 \\
2 & 스파이크 채널의 용량은 $R_{\max}\approx r\log_2\frac{e}{r\Delta t}$~\eqref{eq:spikeentropy}이고, 거의 전부가 스파이크 시각에 있다(명제~\ref{prop:capacity}) & 길이 $\Delta t$ 칸으로 된 이진 문자열의 엔트로피. 장의 수치: 초당 $r=10$, $\Delta t=1\ms$에서 스파이크당 $8$비트, $\Delta t=10\ms$에서 약 $5$비트, 스파이크 계수기는 초당 $2$비트 미만 & \cite{MacKay1952} & 이론 \\
3 & 감각 뉴런은 스파이크당 1비트에서 몇 비트를 전달하며, 가장 좋은 경우 한계의 절반까지 이른다 & 자극을 알고 있는 감각 뉴런의 스파이크로 자극을 복원; 책에 측정 요약이 있다 & \cite{Rieke1997} & 측정됨 \\
4 & 스파이크 발생기는 정확하다. 정확한 시각은 입력의 빠른 변화가 정한다 & 쥐 피질 절편의 뉴런: 빠르게 요동하는 반복 전류에는 스파이크가 $1\ms$보다 나은 정밀도로 반복되고, 일정한 전류에는 시각이 흩어진다 & \cite{Mainen1995} & 측정됨 \\
5 & 시냅스는 믿을 수 없다. 방출의 무작위성 때문에 스파이크의 절반 넘게가 수신자에 이르지 못한다 & 해마 피라미드 뉴런 입력의 최소 자극(절편): 스파이크의 절반 넘게가 응답을 일으키지 않는다. 자극 문턱의 요동, 축삭 가지점에서의 전도 실패, 낮은 실험 온도는 배제했다. 남는 것은 확률적 방출이다 & \cite{AllenStevens1994} & 측정됨 \\
6 & 살아 있는 피질의 스파이크 흐름은 무작위로 보인다. 간격은 푸아송 흐름처럼 흩어지고, 스파이크 수의 분산은 평균에 가깝다. “잡음”의 일부는 동물의 움직임에 관한 메시지이다 & 깨어 있는 원숭이의 시각 영역 V1과 MT 기록: 간격의 변동 계수는 약 $1$이고, 스파이크 수의 분산 대 평균 비는 무작위 과정과 같다. 독립적인 작은 입력을 합하는 모델은 훨씬 작은 산포를 준다. 생쥐 시각 피질에서는 변동의 상당 부분이 동물 자신의 움직임을 찍은 영상으로 예측된다 & \cite{Softky1993,Stringer2019} & 측정됨 \\
7 & 발화율 부호는 느리다. 오차는 $1/\sqrt{rT}$~\eqref{eq:rateerror}인데, 뇌는 $\sim150\ms$ 만에 그림을 알아본다 & 공식은 푸아송 통계로, 장에서 유도했다. 실험: 사람이 $20\ms$ 동안 보여 준 낯선 사진에 동물이 있는지 판단했다. “예”와 “아니오”에 대한 뇌 전위는 약 $150\ms$ 뒤에 갈라진다. 이 시간을 $10$--$20\ms$씩 열 개쯤의 단계로 나눈 것은 장의 추정이지 측정이 아니다 & \cite{Thorpe1996} & 측정됨 \\
8 & 무리에 대한 평균은 상관으로 제한된다. 오차는 기껏해야 $1/\sqrt c$배 줄어든다~\eqref{eq:correlated}(명제~\ref{prop:pooling}) & 원숭이 MT 영역에서 동시에 기록한 뉴런 쌍: 스파이크 수의 상관 계수는 평균 $0.12$이며, 이론적 분석은 이 정도 상관만으로도 무리에서 얻는 이득이 제한됨을 보여 준다. 이론: 한계를 정하는 것은 상관 가운데 신호를 닮은 부분뿐이다. 반대 입장의 모델: 뉴런 $50$--$100$개 무리의 발화율은 스파이크 간격 하나, $10$--$50\ms$ 안에 읽히며, 피질에서 개별 스파이크의 시각은 쓰이지 않는다 & \cite{Zohary1994,MorenoBote2014,ShadlenNewsome1998} & 측정됨 \\
9 & 수는 첫 스파이크 시각~\eqref{eq:latency}, 무리의 스파이크 순서, 국소 리듬의 위상에 적을 수 있다 & 도롱뇽 망막: 짧게 보여 준 그림의 공간 구조가 출력 뉴런 첫 스파이크의 상대 지연에 적히며, 대비와 거의 무관하다. 쥐 해마의 장소 세포: “자기” 장소를 지날 때 스파이크가 세타 리듬($7$--$12$\,Hz) 주기 안에서 점점 일찍 옮겨 가며, 이동 폭은 $100^\circ$에서 $355^\circ$이다. 피질이 스파이크 시각을 얼마나 쓰는지는 열린 문제이다(8행) & \cite{Gollisch2008,OKeefe1993} & 측정됨 \\
10 & 어떤 부호가 읽힐지는 수신자의 시간 상수가 정한다~\eqref{eq:reader}(명제~\ref{prop:reader}). 활성 피질의 뉴런은 동시성 검출기에 더 가깝다 & 식~\eqref{eq:reader}는 장에서 유도했다. in vivo 기록 총설: 활성 피질에서 막 컨덕턴스는 휴지 상태보다 몇 배 크고, 시간 상수는 그만큼 짧다. 피질이 바로 이런 식으로 부호를 전환한다는 것은 직접 보이지 않았다 & \cite{Destexhe2003} & 가설 \\
11 & 고갈 시냅스는 발화율의 변화, 본질적으로 $\Delta r/r$를 전달한다~\eqref{eq:depression}(그림~\ref{fig:depression}) & 피질 피라미드 뉴런의 쌍 기록: 고갈 속도는 방출 확률이 정한다. 고갈이 느리면 응답이 발화율을 반영하고, 빠르면 일치를 반영한다. 쥐 피질에서 측정한 고갈에 기초한 모델: 빠른 입력과 느린 입력에서 같은 상대 발화율 변화가 같은 응답을 준다. 즉 시냅스는 $\Delta r/r$를 전달한다. 수치 $1.7\to2.2$, $6.7$, $90\ms$는 장의 계산이다 & \cite{Tsodyks1997,Abbott1997depression} & 측정됨 \\
12 & 버스트는 “대시”이다. 버스트가 사라질 확률 $(1-p)^k$~\eqref{eq:burst}는 작고, 촉진이 이를 더 줄인다 & 총설: 믿을 수 없는 시냅스에서 단일 스파이크는 자주 사라지지만, 버스트는 촉진 덕분에 믿을 만하게 전달된다. 버스트는 시냅스의 장기 변화를 일으킨다. 뇌가 버스트를 별도의 기호로 읽는다는 것은 가설이다 & \cite{Lisman1997,AllenStevens1994} & 가설 \\
13 & 빠른 \nt{GABA} 뉴런은 리듬과 “구두점”을 정한다. 또 다른 부류는 억제하는 뉴런을 억제하여 관문을 연다(명제~\ref{prop:punctuation}) & 피질 \nt{GABA} 뉴런 부류들의 성질에 관한 총설. 광유전학: 빠른 \nt{GABA} 뉴런을 리듬 있게 흥분시키면 감마 리듬이 생기고, 자극에 대한 응답은 주기의 위상에 의존한다. 생쥐 시각 피질에서 PV 뉴런은 서로를, SST 뉴런은 나머지 모든 부류를, VIP 뉴런은 주로 SST를 억제한다. 깨어 있는 생쥐에서 VIP 뉴런을 흥분시키면 \nt{GLUT} 뉴런의 억제가 풀린다. VIP 뉴런은 강화 신호로 켜진다. 일반 규칙으로서의 역할 분담은 가설이다 & \cite{Tremblay2016,Cardin2009,Pfeffer2013,Pi2013} & 가설 \\
14 & 에너지는 평균 발화율을 초당 스파이크 1개 정도로 제한한다~\eqref{eq:energy} & 해부학적, 생리학적 자료로 본 설치류 회색질의 에너지 예산: 스파이크와 시냅스 전류가 신호 전달 에너지의 $47\,\%$와 $34\,\%$이다. 동시에 활성일 수 있는 뉴런은 $\sim15\,\%$를 넘지 못한다. 사람 피질에서는 동시에 강하게 활성일 수 있는 뉴런이 $\sim1\,\%$를 넘지 못한다. 스파이크당 ATP $\sim10^9$개와 신호 전달 전력 $\sim10$\,W는 이 계산에 기초한 장의 추정이다 & \cite{Attwell2001,Lennie2003} & 이론 \\
15 & 부호는 희소하며 줄당 비트 수가 최대가 되도록 맞추어져 있다. 피질은 정밀도를 에너지와 맞바꾼다(명제~\ref{prop:economy}) & 절약 부호 가설. 근거: 먹이를 제한한 생쥐의 시각 피질에서 AMPA 수용체 컨덕턴스와 시냅스의 ATP 소비가 줄고($-29\,\%$), 스파이크 발화율은 유지되며, 방위 조율은 $32\,\%$ 넓어진다 & \cite{Barlow1961,Padamsey2022} & 가설 \\
\bottomrule
\end{longtable}}

\section{\ref{sec:serocomputer}장. \nt{SERO} 뉴런}

\nt{SERO} 뉴런 자체의 성질(리듬, 자가 억제, 수용체에 따른 부호, 하위 시스템)은 직접 측정되었다. 이 장의 계산, 즉 조절기, 필터, 논리는 이 장의 식에서 나온다. $\tau_c$, 세로토닌의 확산 계수, 방출 지점의 수는 추정이다. 고리가 뇌에서 수준 조절기로 작동한다는 것은 가설이고(봉선핵에서 자가 억제는 점대점이며 짧은 멈춤만 만든다), \nt{SERO}가 지평을 정한다는 것은 이를 지지하는 행동 실험이 있는 가설이다.

{\footnotesize
\setlength{\tabcolsep}{3pt}\renewcommand{\arraystretch}{1.15}
\begin{longtable}{>{\raggedright}p{0.5cm}>{\raggedright}p{4.0cm}>{\raggedright}p{5.6cm}>{\raggedright}p{2.3cm}>{\raggedright\arraybackslash}p{1.7cm}}
\toprule
No. & 명제 & 실험과 측정 내용 & 출처 & 상태 \\
\midrule\endfirsthead
\toprule
No. & 명제 & 실험과 측정 내용 & 출처 & 상태 \\
\midrule\endhead
1 & 세로토닌 작용의 부호는 수신자의 수용체가 고른다(명제~\ref{prop:receptor}) & 세로토닌 수용체는 약리와 기전에 따라 여러 계열로 나뉜다. 5-HT$_{1A}$는 G 단백질을 통해 칼륨 채널을 열어 세포를 억제하고, 5-HT$_{2A}$와 이온성 5-HT$_3$는 흥분시킨다. 한 신경전달물질이 세포에 따라 반대 응답을 낸다. & \cite{BarnesSharp1999} & 측정됨 \\
2 & 깨어 있을 때 \nt{SERO} 뉴런은 초당 몇 개의 스파이크를 고르게 낸다 & 자유롭게 움직이는 고양이의 등쪽 봉선핵 뉴런 기록: 조용히 깨어 있을 때 $2.8\pm0.2$ Hz의 느리고 규칙적인 발화. 서파 수면에서는 발화율이 $34$--$68\,\%$, 렘수면에서는 $98\,\%$ 떨어진다. 마취한 쥐에서는 $1.7\pm0.5$ Hz로 매우 규칙적이다. & \cite{TrulsonJacobs1979, GagnonParent2014, JacobsAzmitia1992} & 측정됨 \\
3 & \nt{SERO} 뉴런은 페이스메이커다. 스파이크마다 밀어 줄 필요는 없지만 고른 배경 입력은 필요하다(절편에서는 $\alpha_1$ 수용체를 통한 노르아드레날린, 생체에서는 \nt{NORA}에서 온다) & 뇌 절편에서 세포는 느리고 고르게 발화하며, 스파이크마다 $200$--$800\ms$의 후과분극이 뒤따른다. 절편에서 자발적으로 발화하는 세포는 생체보다 적다. 고른 리듬에는 $\alpha_1$ 수용체를 통한 노르아드레날린의 지속적 입력이 필요하고, 이를 차단하면 리듬이 꺼진다. & \cite{VandermaelenAghajanian1983} & 측정됨 \\
4 & 수용체는 거의 모두 대사성이다. 응답은 느려서 약 1초 안에 오르고 내린다 & 5-HT$_3$을 뺀 모든 수용체 계열은 G 단백질과 결합한다. 절편에서 유발된 세로토닌 방출은 5-HT$_{1A}$를 통해 1초 안에 오르고 내리는 억제 전류를 만든다. & \cite{BarnesSharp1999, CourtneyFord2016} & 측정됨 \\
5 & 목표 영역에서 세로토닌은 틈뿐 아니라 부피로도 나온다. 봉선핵 자체에서 이웃 억제는 점대점이다 & 절편에서의 고속 주사 순환 전압전류법: 스파이크 하나당 방출량이 시냅스가 있는 영역과 시냅스가 거의 없는 봉선핵에서 같다. 재흡수와 수용체를 차단해도 달라지지 않는다. 말단당 분자 수가 수송체와 수용체보다 훨씬 적고, 세포 밖 농도는 주 수용체의 친화도와 일치한다. 봉선핵에서 5-HT$_{1A}$를 통한 억제는 점대점이다. 수송체가 세로토닌이 틈 밖에 쌓이지 못하게 한다. & \cite{Bunin1998, CourtneyFord2016} & 측정됨 \\
6 & 식~\eqref{eq:seroneuron}의 농도는 재흡수로 줄어든다. $\tau_c$가 1초 정도라는 것은 추정이다 & 생체 뇌에서의 전압전류법: 세포 밖 세로토닌을 제한하는 것은 합성이 아니라 주로 재흡수와 분해다. 인용한 연구에 $\tau_c$의 직접 측정은 없다. 크기 정도는 이 장의 추정이다. & \cite{Hashemi2012serotonin} & 측정됨; $\tau_c$는 추정 \\
7 & 페이스메이커 하나로 만드는 NOT과 NOR; \nt{SERO} 논리의 완전성(명제~\ref{prop:serocomplete}, 그림~\ref{fig:serogates}) & 식~\eqref{eq:seroneuron}에서 나온다. 억제되는 페이스메이커는 NOR이고, NOR은 함수적으로 완전하다. \nt{SERO} 뉴런으로 논리를 만든 실험은 없다. 부피가 분리되어 있어야 한다는 조건은 뇌에서 성립하지 않는다. & 이 장의 유도 & 이론 \\
8 & 세로토닌은 \nt{SERO} 뉴런 자신에 있는 수용체를 통해 그 뉴런을 억제한다 & 마취한 쥐에서 5-HT$_{1A}$ 작용제를 정맥 주사하거나 이온영동으로 주면, 세로토닌 자체와 같은 용량--반응 관계로 등쪽 봉선 뉴런의 발화가 억제된다. 5-HT$_{1B}$ 작용제는 거의 효과가 없다. 절편에서 이웃 세포의 내인성 세로토닌은 5-HT$_{1A}$를 통해 전류와 발화의 짧은 멈춤을 만들지만 평균 발화율은 바꾸지 않는다. & \cite{Sprouse1987, CourtneyFord2016} & 측정됨 \\
9 & 모델에서 공통 부피를 통한 고리는 수준 조절기다. 편차와 시간 상수가 $1+G$배 줄어든다~\eqref{eq:feedback}. 고리가 뇌에서 이렇게 작동한다는 것은 가설이다 & 음의 되먹임 증폭기의 고전적 공식이며, 이 장에서 새로 유도했다. \nt{SERO} 시스템의 고리 이득 $G$는 측정되지 않았다. 측정된 것은 절편에서 자가 억제가 점대점이고, 짧은 멈춤을 만들며, 평균 발화율을 바꾸지 않는다는 점이다. & \cite{Black1934feedback, CourtneyFord2016} & 이론; 뇌에서는 가설 \\
10 & 농도는 발화율의 이동 평균, 즉 저역 통과 필터다~\eqref{eq:lowpass}, 그림~\ref{fig:lowpass} & 선형 방정식~\eqref{eq:seroneuron}의 해. 그림은 $\tau_c=1\,\text{s}$에서 이 공식으로 계산한 것이다. \texttt{models/}에 별도 모델은 없다. & 이 장의 유도 & 이론 \\
11 & 틈의 굴곡은 작은 분자의 확산을 약 $2.6$배 늦춘다 & 점원 방법: 테트라메틸암모늄을 이온영동으로 주입하고 이온 선택성 전극으로 그 농도를 절편과 생체 뇌에서 기록한다. 굴곡도 $\lambda\approx1.6$, 즉 유효 $D$는 자유 확산보다 $\lambda^2\approx2.6$배 작다. 세포 사이 부피의 비율은 약 $0.2$. 이 연구들에서 조직 속 세로토닌 자체의 확산 계수는 측정하지 않았다. 이 장에서 그 값은 추정이다. & \cite{Sykova2008} & 측정됨 \\
12 & 독립 “전선”의 길이 $\ell\sim\sqrt{D\tau}\approx20$~\textmu m~\eqref{eq:diffusionlength}; 전선의 수는 이런 영역의 수로 제한된다(명제~\ref{prop:volumewires}) & 확산 법칙에 $D$와 $\tau$의 추정값(6행과 11행)을 대입한 것이며, 명제~\ref{prop:volumewires}는 그 따름정리다. 체적 전달의 독립 채널 수를 센 직접 실험은 없다. & 이 장의 유도 & 이론 \\
13 & \nt{SERO} 뉴런 하나의 출력은 뇌의 넓은 영역에 퍼져 있다. 방출 지점은 추정으로 $\sim10^5$개 & 쥐 등쪽 봉선 뉴런 하나하나의 완전 재구성: 모든 상행 축삭이 많이 갈라지며, 세포 하나의 축삭 총길이는 최대 $18.7$~cm, 한 세포의 가지가 선조체, 전전두피질, 편도체에 닿는다. 세포당 방출 지점 수는 직접 세지 않았다. & \cite{GagnonParent2014} & 측정됨; $10^5$은 추정 \\
14 & \nt{SERO} 뉴런은 출력과 응답이 다른 몇 개의 하위 시스템으로 나뉜다 & 생쥐에서의 바이러스 유전 표지: 편도체로 가는 뉴런과 전두피질로 가는 뉴런은 분지가 거의 겹치지 않고, 서로 다른 입력을 받으며, 불쾌한 자극에 반대로 반응한다. 각 집단을 켜고 끄면 행동이 서로 다르게 바뀐다. & \cite{Ren2018} & 측정됨 \\
15 & 지수형 할인에서 인내 조건은 $D<\tau_\gamma\ln(R/r_0)$~\eqref{eq:patience}; 측정된 쌍곡선형 할인에서는 $D<\tau_\gamma(R/r_0-1)$ & 할인 $e^{-D/\tau_\gamma}$ 또는 $1/(1+D/\tau_\gamma)$에서 나오며, 이 장에서 유도했다. 몇 초 지연에서 선택하는 원숭이: 할인은 지수보다 쌍곡선에 가깝다. 지연 보상을 알리는 신호에 대한 \nt{DOPA} 뉴런의 응답도 지연에 따라 쌍곡선으로 떨어진다. & 이 장의 유도; \cite{KobayashiSchultz2008} & 이론 \\
16 & \nt{SERO}가 지평 $\tau_\gamma$를 정한다(\ref{sec:horizon}절) & 생쥐에서 등쪽 봉선 \nt{SERO} 뉴런을 광유전학으로 켜면 지연 보상을 더 오래 기다리고, 드물어진 보상을 더 끈질기게 찾는다. 사람에서 세로토닌을 낮추면(트립토판을 뺀 식단) 지연 보상의 할인이 가팔라진다. 농도가 어떻게 $\tau_\gamma$로 바뀌는지는 알려져 있지 않다. & \cite{Doya2002, Miyazaki2014, Lottem2018, Schweighofer2008} & 가설 \\
\bottomrule
\end{longtable}}

\section{\ref{sec:dofa}장. \nt{DOPA}를 이용한 학습}

시냅스의 메커니즘(꾸러미, 동시성 검출기, 칼슘, 가소성 창)과 예측 오차로서의 도파민의 행동은 직접 측정되었다. 규칙들이 기울기로 이어진다는 것과 브로드캐스트의 대가는 정리다. 선택 학습의 결과는 이 책의 모델로 계산한 것이다. 오차를 계산하는 회로는 가설이다.

{\footnotesize
\setlength{\tabcolsep}{3pt}\renewcommand{\arraystretch}{1.15}
\begin{longtable}{>{\raggedright}p{0.5cm}>{\raggedright}p{4.0cm}>{\raggedright}p{5.6cm}>{\raggedright}p{2.3cm}>{\raggedright\arraybackslash}p{1.7cm}}
\toprule
No. & 명제 & 실험과 측정한 것 & 출처 & 상태 \\
\midrule\endfirsthead
\toprule
No. & 명제 & 실험과 측정한 것 & 출처 & 상태 \\
\midrule\endhead
1 & 인공 신경망과 같은 형태의 오차 역전파는 뇌에서 발견되지 않았다 & 직접적인 실험은 없다. 이것은 부재에 관한 명제다. 리뷰들은 이 알고리즘에 무엇이 필요한지(순방향 연결을 되풀이하는 역방향 연결, 별도의 단계)와 어떤 생물학적 근사가 제안되었는지를 다룬다. & \cite{Lillicrap2020, WhittingtonBogacz2019} & 가설 \\
2 & 가중치는 방출 부위 수, 방출 확률, 한 꾸러미에 대한 응답으로 분해된다: $w=N_s\,p\,A_1$~\eqref{eq:quantal} & 방출을 낮춘 개구리 신경근 시냅스: 수신자의 응답은 꾸러미 하나에 대한 자발적 미니 응답의 정수배로 계단식으로 변한다. 스파이크당 꾸러미 수는 푸아송 법칙을 따른다. & \cite{delCastillo1954} & 측정됨 \\
3 & 수신자의 수용체는 동시성 검출기이고, 칼슘이 가중치 변화를 일으킨다 & 생쥐 뉴런의 패치 클램프: 휴지 전위에서는 Mg$^{2+}$ 이온이 글루탐산으로 열린 채널을 막고, 탈분극되면 빠져나간다. 해마 절편: 수신자 안의 칼슘 킬레이터는 장기 강화를 막고, 세포 안에서 빛으로 칼슘을 풀어 주면 그것만으로 전달이 강화된다. & \cite{Nowak1984, Malenka1988calcium} & 측정됨 \\
4 & 결정은 어느 쪽이든 실행한다: 수신자에서는 $A_1$이 커지고, 송신자에서는 $p$가 바뀐다 & 가시 하나 위에서 빛으로 풀어 준 글루탐산은 가시와 그 수용체를 흐르는 전류를 빠르고 지속적으로 키운다. 강화에는 AMPA 수용체의 삽입이 따른다. 수신자의 탈분극은 송신자의 방출을 한동안 억누른다. 이 효과는 틈을 거슬러 가는 엔도카나비노이드가 전달한다(\nt{GABA} 입력에서 확인됨). 단기 고갈은 송신자의 방출 확률에 달려 있다. & \cite{Matsuzaki2004, Malinow2002, WilsonNicoll2001, Tsodyks1997} & 측정됨 \\
5 & 송신자와 수신자가 함께 활동하면 시냅스가 강해진다~\eqref{eq:hebb} & 마취한 토끼: 해마 입력 경로에 짧은 스파이크 열을 준 뒤 응답이 $30$~min에서 $10$~h까지 강화된다. 해마 절편에서 수신자의 스파이크는 같은 순간 활동한 시냅스만 강화한다. & \cite{Bliss1973LTP, Kelso1986hebb} & 측정됨 \\
6 & 규칙~\eqref{eq:oja}은 입력 상관행렬의 주 고유벡터를 찾는다(명제~\ref{prop:oja}) & 정리이며 증명은 본문에 있다. 뉴런이 자기 입력의 주성분을 배운 실험은 없다. 시냅스에서 가중치 증가를 제한하는 메커니즘은 밝혀지지 않았다. & \cite{Oja1982} & 이론 \\
7 & 시각에 따른 헤브 규칙: 창은 약 $\pm20\ms$(그림~\ref{fig:stdp}). 반복되는 순서열에서 사슬이 자란다 & 배양한 해마 뉴런 쌍: 송신자의 스파이크 뒤 $20\ms$ 안의 수신자 스파이크는 지속적 강화를, $20\ms$ 앞의 스파이크는 약화를 준다. 둘 다 NMDA 수용체에 의존한다. 그림의 비율 $A_-=0.6A_+$는 예시다. 이렇게 신경망에서 사슬이 자란다는 것은 직접 측정되지 않았다. & \cite{BiPoo1998} & 측정됨 \\
8 & 적격 흔적~\eqref{eq:threefactor}: 함께 일어난 활동 뒤 $1$--$2\,\text{s}$ 안에 온 도파민은 연결을 강화하고, 더 늦으면 그렇지 않다(명제~\ref{prop:threefactor}) & 선조체 절편, 글루탐산 입력과 도파민 입력의 개별 광유전학 자극: 도파민이 글루탐산 입력 뒤 $0.3$--$2\,\text{s}$ 안에 올 때만 가시가 자란다. 창은 cAMP의 빠른 상승과 분해로 정해진다. 피질에서는 짝짓기가 강화와 약화에 대해 따로따로 흔적을 남기고, 모노아민 수용체가 이것을 몇 초 규모의 창 안에서 장기 변화로 바꾼다. & \cite{Yagishita2014, He2015eligibility} & 측정됨 \\
9 & 몇 초 길이의 창은 도파민 없이도 있다: 세 번째 신호는 수신자의 강한 흥분이다 & 살아 있는 생쥐, CA1 영역: 수신자의 플래토 사건 하나가 그 앞뒤 몇 초 안에 온 입력을 강화하고, 한 번 지나가는 것만으로 장소장을 만든다. & \cite{Bittner2017} & 측정됨 \\
10 & 3요소 규칙의 평균 걸음은 보상의 기울기를 따른다~\eqref{eq:perturbation}(명제~\ref{prop:gradient}) & 무작위 이탈로 학습할 때의 기울기에 관한 정리. 본문에서는 작은 잡음에 대해 유도했다. & \cite{Williams1992} & 이론 \\
11 & 잡음은 탐색이다: 신경망은 자기의 무작위 이탈로부터 배운다 & 어린 금화조: 기저핵 회로의 출력 핵을 끄면 노래의 변동성이 급격하고 가역적으로 사라진다. 성체 금화조: 음높이가 평균의 한쪽에 있는 음절을 부를 때 방해음을 주면 새는 음높이를 빠르게 반대쪽으로 옮긴다. 자기 산포로부터 배우는 것이다. & \cite{Olveczky2005, TumerBrainard2007} & 측정됨 \\
12 & 두 선택지 중 고르기는 $\tau_e=1\,\text{s}$, $\delta=R-\bar R$에서 학습되고, $\tau_e=50\ms$에서는 학습되지 않는다. 기대값을 빼지 않으면 가중치가 한없이 자라고, 학습률이 크면 신경망이 더 나쁜 선택을 굳힐 수 있다 & \texttt{models/dofa\_learning.py}, $\eta=0.05$, 시도 $500$회, 시드 $6$개. $\tau_e=1\,\text{s}$: $A$ 선택 비율 $0.65$--$0.79\to0.96$--$1.00$, 가중치 $0.85$--$1.33$. $\tau_e=50\ms$: $0.38$--$0.55$, 가중치 변화 없음. $\delta=R$: 승자의 가중치 $860$--$1190$. $\delta=R$, $\eta=0.5$, 시드 $3$개: 하나가 $B$에 고정($0.04\to0$). 스크립트가 네 경우를 모두 출력하며, 재계산 결과는 본문과 일치했다. & \texttt{models/\allowbreak dofa\_\allowbreak learning.py} & 모델 \\
13 & 하나의 공통 신호로 학습하는 시간은 잡음을 내는 뉴런 수에 비례해 늘어난다(명제~\ref{prop:broadcastcost}) & 선형 신경망의 학습 곡선: 뉴런 섭동은 정확한 기울기보다 출력 뉴런 수만큼 느리고, 가중치 섭동은 거기에 입력 수만큼 더 느리다. & \cite{Werfel2005} & 이론 \\
14 & \nt{DOPA}의 출력은 무엇보다 행동 선택 영역으로 간다 & 쥐의 흑질선조체 \nt{DOPA} 뉴런 개개의 축삭 완전 재구성: 세포 하나의 가지가 선조체 부피의 $0.45$--$5.7\,\%$(평균 $2.7\,\%$)를 덮고, 선조체 밖에는 가지가 거의 없다. & \cite{Matsuda2009} & 측정됨 \\
15 & 피질은 자기 입력을 예측하도록 배우며, 오차는 그 자리에서 계산된다 & 리뷰: 감각 피질에는 기대한 입력과 실제 입력의 불일치에 대한 응답이 있다. 이것이 피질 학습의 일반 규칙이라는 것은 증명되지 않았다. & \cite{Keller2018} & 가설 \\
16 & 도파민은 오차~\eqref{eq:ctd}처럼 행동한다: 급증, 예고 신호로의 이동, 급감(그림~\ref{fig:ctdcases}) & 원숭이의 \nt{DOPA} 뉴런: 예상 못 한 주스에 급증. 학습 뒤에는 예고 신호에 급증이 생기고 약속된 주스에는 응답이 없다. 주스가 오지 않으면 급감. 연속 형태~\eqref{eq:ctd}는 이론이다. & \cite{Schultz1997, Doya2000} & 측정됨 \\
17 & $\dd V/\dd t$를 계산하고 기대를 빼는 회로 & 생쥐, 배쪽 피개 영역의 기록과 광유전학: \nt{DOPA} 뉴런은 보상에서 기대를 뺀다. 이웃한 \nt{GABA} 뉴런은 보상이 기대될 때 이들을 억제하며, 이 뉴런들을 흥분시키거나 억제하면 오차가 바뀐다. 도함수가 어떻게 계산되는지는 밝혀지지 않았다. & \cite{Eshel2015arithmetic} & 가설 \\
18 & \nt{DOPA} 뉴런은 페이스메이커이며, 급감은 급증보다 얕다 & 절편에서 \nt{DOPA} 뉴런은 느린 페이스메이커 전위 위에서 자발적으로, 매우 규칙적으로 발화한다. 원숭이에서 발화 빈도는 양의 오차를 정량적으로 따르지만 음의 오차는 약하게만 전달한다. & \cite{GraceOnn1989, Bayer2005} & 측정됨 \\
19 & 행위자와 비평자(그림~\ref{fig:actorcritic}) & 알고리즘은 강화 학습 이론에서 왔다. 사람(fMRI)에서 배쪽 선조체는 선택이 있든 없든 비평자처럼 행동하고, 등쪽 선조체는 행동을 골라야 할 때만 그렇다. & \cite{SuttonBarto2018, ODoherty2004actor} & 가설 \\
20 & 두 번째 계열의 수용체를 가진 뉴런에서는 규칙의 $\delta$ 부호가 뒤집힌다 & 선조체 절편: D1 수용체를 가진 뉴런에서는 짝짓기가 D1이 필요한 강화를, D2를 가진 뉴런에서는 D2가 필요한 약화를 준다. & \cite{Shen2008} & 측정됨 \\
\bottomrule
\end{longtable}}

\section{\ref{sec:dofaalt}장. \nt{DOPA}의 다른 이론들}

확장된 그림들은 저마다 도파민의 직접 측정(반전점의 흩어짐, 불쾌한 자극과 새로운 자극에 대한 반응, 완만한 상승, 맛 교체에 대한 반응)에 기대고 있다. 그러나 그것을 설명하는 식들은 이론이며, 그 가운데 두 이론 사이에서는 실험 결과가 아직 서로 모순된다.

{\footnotesize
\setlength{\tabcolsep}{3pt}\renewcommand{\arraystretch}{1.15}
\begin{longtable}{>{\raggedright}p{0.5cm}>{\raggedright}p{4.0cm}>{\raggedright}p{5.6cm}>{\raggedright}p{2.3cm}>{\raggedright\arraybackslash}p{1.7cm}}
\toprule
No. & 명제 & 실험과 측정한 것 & 출처 & 상태 \\
\midrule\endfirsthead
\toprule
No. & 명제 & 실험과 측정한 것 & 출처 & 상태 \\
\midrule\endhead
1 & 비대칭 규칙~\eqref{eq:asymmetric}은 점~\eqref{eq:expectile}으로 수렴한다. $\kappa=1/2$이면 평균이고, 확률 $1/2$의 한 방울이면 $V=\kappa$이다(그림~\ref{fig:expectiles}) & 점~\eqref{eq:expectile}은 익스펙타일, 곧 양의 편차와 음의 편차에 서로 다른 가중치를 준 최소제곱 문제의 해다. 이 장의 예에 대한 유도는 본문에 있다. & \cite{NeweyPowell1987}; 이 장의 유도 & 이론 \\
2 & \nt{DOPA} 뉴런마다 반전점이 다르고, 반전점은 비대칭 순으로 정렬된다(명제~\ref{prop:expectile}) & 생쥐, 배쪽 피개 영역의 \nt{DOPA} 뉴런을 광유전학으로 표지, 크기가 다른 보상: 급증이 급감으로 바뀌는 점이 뉴런마다 다르다. 반응 비대칭이 큰 뉴런일수록 그 점이 높다. 무리 반응으로 보상 분포의 모양이 복원된다. 이것이 흩어짐의 주된 기능이라는 것은 가설이다. & \cite{Dabney2020} & 측정됨 \\
3 & 일부 \nt{DOPA} 뉴런은 불쾌한 자극에 급증으로 반응한다 & 원숭이: 어떤 \nt{DOPA} 뉴런은 보상 예고 신호에 흥분하고 눈에 부는 공기 예고 신호에 억제되며, 다른 뉴런은 둘 다에 흥분한다. 두 무리는 중뇌의 서로 다른 부위에 있다. & \cite{Matsumoto2009, BrombergMartin2010} & 측정됨 \\
4 & 오차의 크기 또는 새로움이 학습률을 정한다 & \nt{DOPA} 뉴런의 중요도 신호가 학습률을 바꾼다는 직접 실험은 인용된 연구에 없다. 리뷰는 “주의, 중요함” 신호의 역할을 제안한다. & \cite{BrombergMartin2010} & 가설 \\
5 & 위험 축을 위한 별도의 전선 & 생쥐: 선조체 꼬리로 가는 \nt{DOPA} 뉴런은 보상이 아니라 새롭고 강한 자극에 반응한다. 이 뉴런을 파괴하면 새로운 물체에 대한 회피가 약해진다. & \cite{Menegas2018} & 측정됨 \\
6 & 최적 행동 시간 $\tau_a^*=\sqrt{C/\bar R}$~\eqref{eq:vigor} & 합 $C/\tau_a+\bar R\tau_a$의 최소. 유도는 본문에 있다. 원래 모델에서 행동 속도는 평균 보상과 같은 시간의 가격이 정한다. & \cite{Niv2007}; 이 장의 유도 & 이론 \\
7 & 느린 도파민 수준은 $\bar R$를 싣고 행동 속도를 정한다(명제~\ref{prop:multiplex}) & 쥐, 측좌핵의 미세투석과 전압전류법: 도파민 수준이 분 단위로 보상 빈도와 행동 속도를 따라간다. 사람에서 L-DOPA는 반응 시간이 평균 보상 빈도에 의존하는 정도를 키운다. 느린 수준이 정확히 $\bar R$라는 것은 증명되지 않았다. & \cite{Hamid2016, Beierholm2013vigor} & 가설 \\
8 & 느린 수준은 두 원천에서 합쳐진다. 방출은 스파이크 발화율이 변하지 않아도 출력 말단에서 달라지지만, 발화율 자체에도 완만한 상승이 있다 & 쥐, 한 과제: 측좌핵의 방출은 변하는 보상 기대를 따라가지만, 식별된 배쪽 피개 영역 \nt{DOPA} 뉴런의 발화율은 그렇지 않다. 가상 복도의 생쥐(10행): 보상에 다가갈 때의 완만한 상승이 식별된 \nt{DOPA} 뉴런의 스파이크에도 있다. & \cite{Mohebi2019, Kim2020} & 측정됨 \\
9 & 동물이 먼 보상에 다가가는 동안 도파민이 완만하게 오른다 & 미로의 쥐, 선조체 전압전류법: 목표에 가까워지면서 도파민 농도가 몇 초에 걸쳐 오르고, 기울기는 거리와 보상 크기에 의존한다. & \cite{Howe2013ramp, Hamid2016} & 측정됨 \\
10 & 상승은 추정이 정확해질 때의 오차 $\delta$다~\eqref{eq:ramp}. 복도에서의 도약은 급증을 준다(그림~\ref{fig:ramp}) & 식과 수치 초당 $\delta=0.75\,V$는 본문의 유도. 실험: 가상 복도의 생쥐를 목표 쪽으로 순간 이동. 세포체의 스파이크, 세포체와 출력 말단의 칼슘, 선조체의 농도가 기대 $V$가 아니라 오차처럼 반응한다. 두 설명 중 어느 것이 모든 출력에서 옳은지는 논쟁 중이다. & \cite{Kim2020} & 측정됨 \\
11 & 뒤돌아보기: 도파민은 인과 연관의 척도~\eqref{eq:retro}를 알린다 & 생쥐, 이 이론과 오차~\eqref{eq:ctd}가 서로 다른 것을 예측하는 실험에서의 측좌핵 도파민 방출: 여러 실험에서 방출이 앞의 이론을 따랐다. & \cite{Jeong2022} & 가설 \\
12 & 학습 중 도파민 반응은 보상에서 예고 신호로 점차 옮겨 간다 & 생쥐, 학습 과정 동안 배쪽 선조체에서 \nt{DOPA} 뉴런 출력 말단의 신호: 급증이 보상 순간에서 예고 신호 순간으로 점차 옮겨 간다. \eqref{eq:ctd}에 의한 학습이 요구하는 그대로다. & \cite{Amo2022} & 측정됨 \\
13 & \nt{DOPA} 뉴런은 가치가 같아도 보상의 맛이 바뀌면 반응한다 & 쥐, 배쪽 피개 영역 뉴런 기록: 주스를 가치가 같은 다른 주스로 바꾸면 오차~\eqref{eq:ctd}가 0인데도 반응이 생긴다. & \cite{Takahashi2017} & 측정됨 \\
14 & 인위적 급증은 두 중립 자극 사이의 연관을 가르친다 & 쥐, 보상 없이 두 자극을 미리 짝짓는 실험: 첫째 자극에서 둘째 자극으로 넘어갈 때 \nt{DOPA} 뉴런을 광유전학으로 켜면 연관이 생기고, 끄면 학습이 방해된다. & \cite{Sharpe2017} & 측정됨 \\
15 & 특징별 오차 벡터~\eqref{eq:vectortd} & 특징 예측 오차 이론. 벡터 형태에 대한 직접 검증은 없다. & \cite{Gardner2018} & 가설 \\
16 & 한 과제에서 \nt{DOPA} 뉴런마다 서로 다른 양을 싣는다 & 가상 복도의 생쥐, \nt{DOPA} 뉴런의 2광자 영상: 어떤 뉴런은 보상과, 어떤 뉴런은 속도와 방향 전환과, 또 어떤 뉴런은 예고 신호와 선택의 정확도와 연관된다. & \cite{Engelhard2019} & 측정됨 \\
17 & 전선 대신 다발(명제~\ref{prop:bundle}) & 2--16행의 요약: 각 분할은 따로따로 측정되었다. 이것들이 모두 한 신호의 부분이라는 통합된 그림은 실험으로 확인되지 않았다. & --- & 가설 \\
\bottomrule
\end{longtable}}

\section{\ref{sec:nora}장. \nt{NORA}}

\nt{NORA} 시스템의 구조, 두 작동 모드, 동공과의 연관(\nt{NORA}만을 통한 것은 아니다), 신호 대 배경 비의 증가, 행동에서의 뒤집힌 U는 직접 측정되었다. 곱셈 이득 $g$는 모델이다. 이득이 선택 모드를 전환하고 역온도로 작동한다는 것은 이 장의 유도 결과다. 이득 조정의 이점은 이 책의 모델로 계산한 것이다. “\nt{NORA}는 탐색의 노브다”와 “\nt{NORA}는 인터럽트다”는 가설이다.

{\footnotesize
\setlength{\tabcolsep}{3pt}\renewcommand{\arraystretch}{1.15}
\begin{longtable}{>{\raggedright}p{0.5cm}>{\raggedright}p{4.0cm}>{\raggedright}p{5.6cm}>{\raggedright}p{2.3cm}>{\raggedright\arraybackslash}p{1.7cm}}
\toprule
No. & 명제 & 실험과 측정한 것 & 출처 & 상태 \\
\midrule\endfirsthead
\toprule
No. & 명제 & 실험과 측정한 것 & 출처 & 상태 \\
\midrule\endhead
1 & 사람의 \nt{NORA} 뉴런은 수만 개이며, 거의 모두 한 무리에 있다 & 사람 뇌줄기의 연속 절편, 티로신 수산화효소 염색: 청반에 $53\,900$개, 이웃한 하위핵에 $6\,260$개의 세포. & \cite{Baker1989LC} & 측정됨 \\
2 & 출력은 거의 뇌 전체로 퍼지지만, 무리의 부분들이 서로 다른 영역을 어느 정도 나누어 맡는다 & 생쥐에서의 바이러스 유전 표지: 청반의 \nt{NORA} 뉴런은 뇌의 넓은 영역에서 입력을 받고, 뇌와 척수의 넓은 영역으로 출력을 보낸다. 쥐: 편도체로 가는 무리는 공포 학습을 강화하고, 전전두피질로 가는 별도의 무리는 공포 소거를 강화한다. & \cite{Schwarz2015LC, Uematsu2017LC, Poe2020} & 측정됨 \\
3 & 배경 모드: 초당 1회 미만에서 몇 회 사이의 고른 스파이크, 발화율은 각성 수준에 따라 변한다 & 쥐, 수면--각성 주기 동안의 청반 뉴런 기록: 발화율은 깨어 있을 때 가장 높고, 서파 수면에서 낮으며, 렘수면에서는 거의 0이다. 발화율 변화가 상태 전환보다 앞선다. & \cite{AstonJonesBloom1981} & 측정됨 \\
4 & 돌발 모드: 과제에 중요한 사건에 대략 결정 순간에 나오는 짧은 버스트 & 원숭이, 드문 표적 찾기 과제: 모든 청반 뉴런이 표적에만 버스트로 응답한다. 표적 뒤 $\sim90\ms$, 손 반응 $\sim200\ms$ 전이다. 수행이 나쁠 때 버스트가 약하다. 양자택일 과제에서 버스트는 자극보다 반응에 더 정확히 맞물리며, 정답과 오답 모두에서 나타난다. & \cite{AstonJones1994vigilance, Clayton2004LC} & 측정됨 \\
5 & 동공 지름은 \nt{NORA}의 부정확한 측정점이다. \nt{NORA}뿐 아니라 \nt{ACET}와 다른 영역도 따른다 & 원숭이: 청반의 스파이크는 세포 하나의 개별 스파이크까지 동공 확장을 예측한다. 비슷하지만 더 늦은 연관이 둔덕과 띠피질에도 있다. 생쥐: 동공 변동은 피질로 가는 노르아드레날린성 출력과 콜린성 출력의 활동을 모두 따른다. & \cite{Joshi2016, Reimer2016} & 측정됨 \\
6 & 노르아드레날린은 배경 활동을 유발 활동보다 더 강하게 억제한다. 곱셈 이득~\eqref{eq:gain}은 이 효과를 전달하는 모델이다 & 깨어 있는 원숭이, 청각 피질, 노르아드레날린 이온영동: 뉴런의 배경 활동이 같은 종의 소리에 대한 응답보다 더 강하게 억제되어 신호 대 잡음비가 커진다. 곱셈 시그모이드~\eqref{eq:gain}는 네트워크 모델에서 가져왔다. 기울기가 문자 그대로 곱해진다는 것은 측정되지 않았다. & \cite{Foote1975NE, ServanSchreiber1990} & 측정됨; $g$는 모델 \\
7 & 이득은 증폭기 뒤의 잡음에 대해서만 $d'$를 높인다~\eqref{eq:gainsnr}(명제~\ref{prop:gainnoise}) & 이 장의 유도. 네트워크 모델에서 이득은 신호 검출을 개선한다. 증폭기 앞과 뒤의 잡음을 분리해 이 차이를 검증한 실험은 없다. & 이 장의 유도; \cite{ServanSchreiber1990} & 이론 \\
8 & 경계 $g\beta=1$은 선택 네트워크를 “숙고 중”에서 “결정함”으로 전환한다~\eqref{eq:wtagain}(명제~\ref{prop:gainwta}, 그림~\ref{fig:pitchfork}) & 서로 억제하는 두 무리의 선형 분석; 이 장의 유도. 뇌에서 이 문턱값을 직접 측정한 것은 없다. & 이 장의 유도 & 이론 \\
9 & 이득은 선택의 역온도다~\eqref{eq:softmax} & 증폭기 뒤의 잡음이 로지스틱이면 선택 확률은 $T=\sigma/g$인 softmax다. 이 장의 유도. & 이 장의 유도 & 이론 \\
10 & \nt{NORA}는 얼마나 탐색할지 정한다. 높은 배경 활동은 작은 유효 $g$(탐색), 결정 순간의 돌발 버스트는 큰 $g$(결정) & 사람: 무작위 선택 직전의 기저 동공은 알려진 최선을 고르기 직전보다 넓다. 동공이 넓은 사람일수록 탐색 성향이 높다. 쥐: 무작위 선택 모드로의 전환은 앞띠피질로 가는 \nt{NORA} 입력이 조절한다. 버스트가 유효 $g$를 높인다는 것은 직접 측정되지 않았다. & \cite{JepmaNieuwenhuis2011, Tervo2014stochastic, AstonJones2005} & 가설 \\
11 & 보상은 $g$에 대해 뒤집힌 U로 의존한다. 빠르게 변하는 세계에서 최선의 $g$가 더 작다(명제~\ref{prop:explore}, 그림~\ref{fig:invertedU}) & \texttt{models/nora\_explore.py}, $4000$번 시도의 실행 $60$회. $1000$번에 한 번 변화: 최선 $g=16$, 비율 $0.976$; $20$번에 한 번: $g=8$, 비율 $0.886$; $g=0.5$에서 $0.77$. 재계산 결과가 본문과 일치했다. & \texttt{models/\allowbreak nora\_\allowbreak explore.py} & 모델 \\
12 & 행동에서의 뒤집힌 U: 배경 활동이 적당하고 버스트가 뚜렷할 때 가장 잘한다 & 원숭이, 4행과 같은 과제: 수행이 좋은 시기에는 배경 발화율이 적당하고 표적에 대한 버스트가 뚜렷하다. 배경 발화율이 높으면 버스트가 없고 오경보가 많다. 배경 활동과 유발 활동의 변화는 수행의 변동을 따라간다. & \cite{Usher1999LC, AstonJones2005} & 측정됨 \\
13 & \nt{NORA}는 예상치 못한 불확실성을, \nt{ACET}는 예상된 불확실성을 알린다 & 두 종류의 불확실성을 다루는 베이즈 추론 이론. 인용한 연구에는 이 신호들을 신경전달물질별로 분리한 직접 실험이 없다. & \cite{YuDayan2005} & 가설 \\
14 & 급격한 변화 뒤 사람은 더 빨리 학습하고 동공이 커진다 & 사람, 갑작스러운 변화가 있는 예측 과제: 동공의 짧은 변화는 변화 확률의 추정을 따르며, 기저 동공과 함께 새 데이터가 추정을 얼마나 바꾸는지(학습률)를 예측한다. 빛이나 과제와 무관한 동공 변화도 이 학습률을 바꾼다. 여기서 동공이 바로 \nt{NORA}를 보여 준다는 것은 5행의 간접 자료에 기댄 추론이다. & \cite{Nassar2012} & 측정됨 \\
15 & 돌발 버스트는 인터럽트로 작동한다. 네트워크가 상태를 초기화한다 & \nt{NORA} 버스트가 네트워크 상태를 초기화한 직접 실험은 없다. 측정된 것은 중요한 사건에 대한 버스트뿐이다(4행). & \cite{BouretSara2005} & 가설 \\
16 & 뜻밖의 일에 따라 $g$나 학습률을 조정해도 최선의 고정 설정보다 거의 낫지 않다 & \texttt{models/nora\_explore.py}, 시기당 $1000$번 시도인 세계($1000$번에 한 번, $20$번에 한 번 변화). 최선의 고정 $g=8$: $0.925$; $g$ 조정: 최대 $0.927$; 학습률 조정: 최대 $0.926$; 고정 학습률 $0.4$: $0.928$. 재계산 결과가 본문과 일치했다. & \texttt{models/\allowbreak nora\_\allowbreak explore.py} & 모델 \\
\bottomrule
\end{longtable}}

\section{\ref{sec:acet}장. \nt{ACET}를 더하다}

아세틸콜린의 세 가지 작용은 절편과 생체 뇌에서 측정되었고, 쓰기와 읽기의 충돌은 이 책의 모델에서 보였다. \nt{ACET}가 쓰기 허용 선으로 쓰이며 예상된 불확실성을 알린다는 것은 실험의 부분적 뒷받침이 있는 가설이다.

{\footnotesize
\setlength{\tabcolsep}{3pt}\renewcommand{\arraystretch}{1.15}
\begin{longtable}{>{\raggedright}p{0.5cm}>{\raggedright}p{4.0cm}>{\raggedright}p{5.6cm}>{\raggedright}p{2.3cm}>{\raggedright\arraybackslash}p{1.7cm}}
\toprule
No. & 명제 & 실험과 측정 내용 & 출처 & 상태 \\
\midrule\endfirsthead
\toprule
No. & 명제 & 실험과 측정 내용 & 출처 & 상태 \\
\midrule\endhead
1 & 뇌의 \nt{ACET} 뉴런은 많지 않고 몇 개의 집단에 모여 있으며, 출력은 피질 전체로 퍼진다. 브로드캐스트 채널이다 & 쥐에서 아세틸콜린 합성 효소에 대한 항체 염색과 역행성 표지: 피질, 해마, 편도체, 후각 망울, 시상은 아세틸콜린을 주로 기저 전뇌와 뇌간의 조밀한 세포 집단 여섯 개에서 받는다 & \cite{Mesulam1983} & 측정됨 \\
2 & 피질의 아세틸콜린 수준은 깨어 있을 때 높고 서파 수면에서 낮다 & 자유롭게 움직이는 고양이의 피질과 해마에서 뇌파 기록과 함께 한 미세투석: 아세틸콜린 방출은 서파 수면보다 조용히 깨어 있을 때 $\sim100\%$, 활발히 깨어 있을 때 $\sim175\%$ 높다. 렘수면에서 피질의 방출은 깨어 있을 때와 같다 & \cite{Marrosu1995,Hasselmo1999} & 측정됨 \\
3 & 신호의 의미는 수용체가 정한다. 아세틸콜린에는 빠른 채널형 수용체와 세포 안 반응을 일으키는 느린 수용체가 있다(명제~\ref{prop:receptor}) & 측정 결과의 총설: 아세틸콜린이 흥분성, 전달, 가소성에 미치는 작용은 방출 장소, 수용체 아형(니코틴성은 이온 채널, 무스카린성은 G 단백질 경유), 수신 세포에 달려 있다 & \cite{Picciotto2012} & 측정됨 \\
4 & 첫째 작용: 외부 입력은 거의 건드리지 않고 내부 연결을 약하게 한다(식~\eqref{eq:acetnet}의 인수 $1-a$) & 쥐 후각 피질 절편: 아세틸콜린 작용제 $100$\,\textmu M에서 같은 세포의 내부 연결(Ib층) 전위는 $60\%$ 넘게, 외부 입력(Ia층) 전위는 $15\%$ 미만으로 줄어든다. 억제는 시냅스 전 작용이다. 해마 절편(CA1): 내부 경로 $89\%$ 대 입력 경로 $40\%$ & \cite{HasselmoBower1992,HasselmoSchnell1994} & 측정됨 \\
5 & 둘째 작용: 입력을 강하게 한다(인수 $\frac{1+a}{2}$) & 신피질 절편: 니코틴성 수용체는 시상에서 오는 입력 시냅스를 강화하고 피질 내부 시냅스에는 작용하지 않는다. 무스카린성 수용체는 두 종류를 모두 약하게 한다. 아세틸콜린은 뉴런의 흥분성을 높인다(총설) & \cite{Gil1997,Hasselmo2006} & 측정됨 \\
6 & 셋째 작용: 가중치 변화를 쉽게 한다(학습률 $\eta a$) & 쥐: 피질 아세틸콜린의 공급원인 기저핵을 전기 자극하면서 소리를 함께 들려주면, 다 자란 동물에서도 청각 피질 지도가 들려준 소리 쪽으로 크게 재편된다 & \cite{KilgardMerzenich1998,Hasselmo2006} & 측정됨 \\
7 & 식~\eqref{eq:acetnet}의 두 번째 식, 즉 드문 패턴용 헤브 규칙은 일부로부터 패턴을 완성하는 연상 메모리를 준다 & 성긴 패턴과 공분산 규칙을 쓰는 네트워크의 용량 계산: 활성 뉴런 비율 $p$가 작으면 네트워크는 $p=1/2$일 때보다 눈에 띄게 많은 패턴을 저장한다 & \cite{Tsodyks1988} & 이론 \\
8 & 낮은 $a$에서의 쓰기는 메모리를 망가뜨린다. 네트워크는 본 것이 아니라 떠올린 것을 기록한다. $a=0.1$에서도 손상은 $a=0$보다 약하지 않다 & \texttt{models/\allowbreak acet\_memory.py}: 뉴런 $200$개, $20$개짜리 패턴 다섯 개, 새 패턴은 그중 하나와 뉴런 $10$개를 공유, 실행 $10$회. $a=0$에서 새 패턴은 기억되지 않고, 옛 패턴은 남의 뉴런 $10$개 중 $5.9$개를 끌고 나온다. $a=0.1$에서 새 패턴은 떠올려지지만($0.97$) 둘이 합쳐진다. 새 패턴은 남의 뉴런 $7.3$개, 옛 패턴은 $7.9$개를 끌고 나온다. $a=0.2$부터는 하나 미만 & 이 장의 유도 & 모델 \\
9 & 명제~\ref{prop:writeread}: 쓰기와 읽기가 똑같이 잘 되는 수준 $a$는 없다(그림~\ref{fig:acetsim}) & 같은 스크립트, 학습률 $\eta a$: 새 패턴은 한 번 제시로 $a\ge0.9$에서 완전히, $a=0.75$에서 $0.8$만큼 기억되고, $a\le0.5$에서는 기억되지 않는다. 절반으로부터의 읽기: $a\le0.2$에서 $1.0$, $a=0.3$에서 $0.96$, $a=0.4$에서 $0.04$ & 이 장의 유도 & 모델 \\
10 & 뇌에서는 브로드캐스트 전선 하나, 즉 \nt{ACET}가 쓰기와 읽기를 전환한다 & 이 발상은 절편 측정을 바탕으로 만든 모델에서 나왔다. 가장 직접적인 실험은 사람에서다. 무스카린성 수용체 차단제 스코폴라민은 새 단어 쌍, 특히 옛 쌍과 겹치는 쌍의 학습을 해치지만, 이미 배운 쌍을 떠올리는 것은 해치지 않는다. 네트워크의 두 모드를 직접 측정한 결과는 없다 & \cite{HasselmoAndersonBower1992,Atri2004} & 가설 \\
11 & 필터~\eqref{eq:kalman}는 최적이다. 입력 가중치와 학습률은 같은 양 $P/R$이고, 정상 상태에서 $P=\sqrt{qR}$, 기억은 $\sqrt{R/q}$이다(명제~\ref{prop:kalman}) & 실험이 필요 없다. 칼만--뷰시 필터의 최적성은 추정 이론의 고전적 결과이고, 정상 상태는 이 장에서 유도했다 & \cite{KalmanBucy1961} & 이론 \\
12 & \nt{ACET}는 $P$와 비슷한 양, 즉 예상된 불확실성을 네트워크에 알린다 & 주의의 확률 모델에서 제안되었다. 아세틸콜린 수준이 추정의 불확실성을 따라간다는 직접 측정은 없다 & \cite{YuDayan2005} & 가설 \\
13 & \nt{ACET} 뉴런의 일부는 보상과 처벌에 20밀리초 남짓 만에, 강화가 뜻밖일수록 더 강하게 응답한다 & 생쥐, 광유전학으로 식별한 기저 전뇌 콜린성 뉴런: 보상과 처벌에 대한 응답은 $18\pm3$\,ms 뒤, 크기는 강화가 뜻밖일수록 커지며, 두 핵의 뉴런에서 응답이 비슷하다 & \cite{Hangya2015} & 측정됨 \\
14 & \nt{NORA}는 세계의 예상치 못한 변화를 알아채고, \nt{ACET}는 네트워크를 쓰기 모드로 유지한다 & 이 분업은 모델에서 제안되었다. 간접 증거: 사람에서 \nt{NORA}와 연관된 동공 지름이 변화와 학습률을 따라간다. \nt{NORA}--\nt{ACET} 쌍을 직접 검증한 실험은 없다 & \cite{YuDayan2005,Nassar2012} & 가설 \\
15 & 서파 수면에서 읽기 모드의 네트워크는 낮에 기록한 것을 재생하고, 이 재생이 그것을 장기 기억으로 옮겨 쓴다 & 쥐 해마 뉴런 앙상블 기록: 낮에 함께 작동한 세포는 이어지는 서파 수면에서 더 자주 함께, 같은 순서로 발화한다(측정됨). 옮겨 쓰기에 관해서는, 학습 뒤 해마 잔물결을 억제하면 공간 기억이 나빠지고 자발적 잔물결을 늘이면 좋아진다. 그 원인이 낮은 \nt{ACET} 수준이라는 것은 증명되지 않았다 & \cite{WilsonMcNaughton1994,Skaggs1996,Girardeau2009,FernandezRuiz2019,Hasselmo1999} & 가설 \\
\bottomrule
\end{longtable}}

\section{\ref{sec:synthesis}장. 기계 전체}

요약 장이므로 새 실험을 소개하지 않는다. 각 행은 그 명제를 다룬 장과 같은 출처에 기댄다. \nt{GLUT} 버스와 \nt{GABA}를 통한 흐름 제어는 확고하게 확립되어 있고, 브로드캐스트 채널의 역할과 그 공동 동작은 대부분 가설이다.

{\footnotesize
\setlength{\tabcolsep}{3pt}\renewcommand{\arraystretch}{1.15}
\begin{longtable}{>{\raggedright}p{0.5cm}>{\raggedright}p{4.0cm}>{\raggedright}p{5.6cm}>{\raggedright}p{2.3cm}>{\raggedright\arraybackslash}p{1.7cm}}
\toprule
No. & 명제 & 실험과 측정한 것 & 출처 & 상태 \\
\midrule\endfirsthead
\toprule
No. & 명제 & 실험과 측정한 것 & 출처 & 상태 \\
\midrule\endhead
1 & 뉴런 $8.6\cdot10^{10}$개에 제어 신호는 네 종류뿐이다~\eqref{eq:machine} & 뇌 균질액에서 세포핵 계수(등방성 분획법): 성인 남성의 뉴런은 $86.1\pm8.1$\,$\times10^9$개 & \cite{Azevedo2009} & 측정됨 \\
2 & 명제~\ref{prop:architecture}의 앞 두 층: 데이터는 \nt{GLUT} 주소 버스를 따라 흐르고, 연결의 부호는 송신자가 정하며, 흐름은 \nt{GABA} 뉴런이 이끌고 정규화한다(\ref{sec:neurontypes}, \ref{sec:gaba}장) & 뉴런당 주된 신경전달물질은 하나(측정 총설); 피드포워드 억제는 피라미드 뉴런이 입력을 합산하는 창을 좁힌다; 전체 활동에 의한 나눗셈은 여러 영역에서 측정됨 & \cite{Strata1999,Pouille2001,Carandini2012} & 측정됨 \\
3 & 소자는 신뢰할 수 없다: 시냅스는 스파이크 대부분을 잃는다(표~\ref{tab:computer}, \ref{sec:code}장) & 해마 절편, 최소 자극: 스파이크의 절반 넘게 CA1에서 응답을 일으키지 못한다; 문턱값과 축삭 전도의 실패는 배제되었고, 원인은 확률적 신경전달물질 방출 & \cite{AllenStevens1994} & 측정됨 \\
4 & 뇌는 $\sim20$와트로 동작하며, 뉴런당 평균 초당 약 한 번 스파이크한다(\ref{sec:code}장); 엔지니어는 아직 이런 기계를 만들지 못했다 & 측정된 비용에서 얻은 추정~\eqref{eq:energy}: 시냅스 작업을 포함해 스파이크당 ATP 분자 약 $10^9$개(설치류 피질); 피질에 대한 상세 추정은 더 낮은 발화율을 준다. 기술 현황은 총설에 따름 & \cite{Attwell2001,Lennie2003,MehonicKenyon2022} & 이론 \\
5 & 대부분 시냅스의 학습은 국소적 적격 흔적과 하나의 공통 수 $\delta$의 곱이다(\eqref{eq:machine}의 둘째 식, \ref{sec:dofa}장) & 원숭이의 \nt{DOPA} 뉴런은 보상 예측 오차에 응답한다; 선조체 절편에서 도파민은 글루탐산 입력 $0.3$--$2$\,s 뒤에 올 때만 가시를 키운다. 적격 흔적이 측정됨 & \cite{Schultz1997,Yagishita2014} & 측정됨 \\
6 & 출력마다 따로 가는 오차 전선은 운동 학습 회로에만 있다(\ref{sec:motorlearning}장) & 소뇌: 등반섬유 신호가 입력과 동시에 오면 그 입력이 약해진다; 원숭이에서 푸르키네 세포의 복합 스파이크는 운동 오차의 방향을 나르고 보정을 일으킨다 & \cite{Ito2001,Herzfeld2018} & 측정됨 \\
7 & 각 브로드캐스트 채널은 몇 개 선로의 다발이다(명제~\ref{prop:bundle}) & \nt{DOPA}: 같은 과제에서 서로 다른 뉴런이 보상, 운동, 자극 특징을 나른다; 빠른 오차와 느린 도파민 수준은 서로 다르다. \nt{SERO}, \nt{NORA}, \nt{ACET}에서는 이런 분해가 덜 연구됨 & \cite{Engelhard2019,Hamid2016,Mohebi2019} & 측정됨 \\
8 & \nt{SERO}: 수준 조절기이자, 가설상 지평 $\tau_\gamma$(\ref{sec:serocomputer}장) & 자기 억제: 자기 5-HT$_{1A}$ 수용체의 작용제가 솔기핵 세로토닌 뉴런을 억제한다. 측정됨. 지평은 이론에서 제안됨; 가장 강한 실험은 생쥐에서 이 뉴런의 광유전학적 활성화가 지연 보상의 기다림을 늘린다는 것 & \cite{Sprouse1987,Doya2002,Miyazaki2014} & 가설 \\
9 & \nt{NORA}: 이득 $g$가 탐색과 결정 사이에서 고른다(\ref{sec:nora}장) & 신호 대 잡음비 증가: 깨어 있는 원숭이의 청각 피질에서 노르아드레날린은 동종의 울음소리에 대한 응답보다 배경 활동을 더 강하게 억제한다. 측정됨; 곱셈적 이득은 모델; 탐색과 결정 사이의 선택에서의 역할은 이론 & \cite{Foote1975NE,ServanSchreiber1990,AstonJones2005} & 가설 \\
10 & \nt{ACET}: 쓰기와 읽기를 전환한다(\ref{sec:acet}장) & 세 작용(내부 연결 약화, 입력 강화, 가소성 촉진)은 측정됨; 모드 전환은 사람에서 스코폴라민 실험으로 간접 지지됨 & \cite{HasselmoBower1992,Gil1997,Atri2004} & 가설 \\
11 & 식~\eqref{eq:machine}은 기계 전체를 기술한다 & \ref{sec:glutcomputer}--\ref{sec:acet}장 모델의 요약이며, 각 부분은 자기 장에 기댄다; 전체로서는 실험으로 검증되지 않음 & 장 안의 유도 & 가설 \\
12 & 노브들은 공통 제어 없이 협조해 동작한다: 오차, 놀람, 쓰기, 복귀(그림~\ref{fig:timeline}) & 실험 없음: 정성적 도식. 개별 고리는 \nt{NORA}와 \nt{ACET}의 분업 이론, 그리고 사람에서 동공과 학습 속도의 관계 & \cite{YuDayan2005,Nassar2012} & 가설 \\
13 & 하나의 공통 신호에 의한 학습은 결과에 영향을 주는 뉴런이 많을수록 느리다(명제~\ref{prop:broadcastcost}) & 선형 네트워크 학습 곡선 계산: 출력 섭동에 의한 학습은 경사 하강법보다 출력 수만큼 느리다 & \cite{Werfel2005} & 이론 \\
14 & 피질이 다층 회로를 어떻게 조정하는지는 알려져 있지 않다; 후보는 스파이크 버스트, 뉴런 가지 안의 계산, 예측에 의한 자기 지도 학습 & 오차 운반자로서의 버스트는 모델; 피질 뉴런 상세 모델의 응답은 $5$--$8$층 네트워크만 재현했다. 모델; 사람 피질 뉴런 가지에서 XOR을 내는 칼슘 스파이크는 절편에서 측정됨; 자기 지도 학습은 증거 총설 & \cite{Payeur2021,Beniaguev2021,Gidon2020,Keller2018} & 가설 \\
\bottomrule
\end{longtable}}

\section{\ref{sec:sensors}장. 기계의 입력}

센서의 구조는 단일 세포 전류 기록, 달팽이관 진동 기록, 망막 세포 계수로 직접 측정되었다. 감도의 한계와 산탄 잡음 식은 물리에서 나온다. 센서의 부호가 최대 정보에 맞춰져 있다는 것은 강한 근거를 가진 가설이다.

{\footnotesize
\setlength{\tabcolsep}{3pt}\renewcommand{\arraystretch}{1.15}
\begin{longtable}{>{\raggedright}p{0.5cm}>{\raggedright}p{4.0cm}>{\raggedright}p{5.6cm}>{\raggedright}p{2.3cm}>{\raggedright\arraybackslash}p{1.7cm}}
\toprule
No. & 명제 & 실험과 측정한 것 & 출처 & 상태 \\
\midrule\endfirsthead
\toprule
No. & 명제 & 실험과 측정한 것 & 출처 & 상태 \\
\midrule\endhead
1 & 센서는 변환기이다: 수용체가 전류를 만들고, 눈과 귀 감각 세포의 출력은 \nt{GLUT} 버스로 가는 글루탐산이다; 첫 세포들은 스파이크를 내지 않는다(그림~\ref{fig:transducer}) & 거북의 막대세포와 원뿔세포는 흥분성 아미노산을 방출한다: 떼어 낸 막 조각의 글루탐산 채널이 이를 잡는다. 쥐 달팽이관의 속털세포: 신경섬유 말단의 전류는 글루탐산 AMPA 수용체를 거치며, 방출 빈도는 세포의 완만한 탈분극과 함께 $150$\,Hz까지 커진다 & \cite{CopenhagenJahr1989,GlowatzkiFuchs2002} & 측정됨 \\
2 & 어둠 속 빛 센서는 광자 하나에 약 1피코암페어의 전류로 응답한다 & 두꺼비 막대세포 바깥 분절에 흡입 전극: 희미한 섬광에 대한 응답이 양자화되어 있고, 사건은 $\approx1$\,pA(암전류의 $5\%$), 퍼짐 $0.2$\,pA, 효율 $0.5$. 사람은 광자 하나가 들어온 것을 우연보다 자주 알아맞힌다 & \cite{Baylor1979,Tinsley2016} & 측정됨 \\
3 & 소리 센서는 1나노미터보다 작은 변위를 알아챈다 & 뫼스바우어 방법, 기니피그 달팽이관 기저막의 $16$--$19$\,kHz 위치: 청신경 응답 문턱에서 막의 속도는 약 $0.04$\,mm/s, 곧 변위 $\approx0.35$\,nm(우리의 환산). 잰 것은 센서 다발이 아니라 막이다 & \cite{Sellick1982,Hudspeth1989} & 측정됨 \\
4 & 어둠 속 정밀도는 광자 산탄 잡음~\eqref{eq:photonnoise}이 제한한다; 약한 빛에서는 적분이 길어진다 & 식은 푸아송 통계에서 나온다. 막대세포에서 희미한 빛의 잡음은 독립적인 양자 사건들의 합과 일치한다; 밝은 배경에서 섬광 응답은 더 작고 짧다. “수십분의 1초”라는 수치는 여기서 따로 검증되지 않음 & \cite{Baylor1979} & 이론 \\
5 & 입력 범위는 빛이 열 자릿수, 소리가 열두 자릿수이고, 스파이크 발화율은 세 자릿수가 안 된다 & 소리: 특성 주파수에서 기저막 응답은 $40$--$80$\,dB 대역에서 기울기 $0.2$\,dB/dB까지로 커지며, 압축은 $100$\,dB 위에서도 보인다. 자릿수 자체는 표준 추정치로, 본문에서는 출처 없이 제시 & \cite{Ruggero1997} & 측정됨 \\
6 & 센서 응답은~\eqref{eq:nakarushton}이며, $\sigma$는 평균 밝기를 따라가 곡선을 로그 축을 따라 이동시킨다(그림~\ref{fig:adaptation}) & 이 식은 처음 물고기 망막의 S 전위에 맞춰졌다. 거북 원뿔세포: 응답은 $V=I V_m/(I+\sigma)$를 따르고 밝기 $\sim3.5$자릿수를 담는다; 배경에 $2$--$3$분 노출하면 곡선은 폭을 유지한 채 수평으로 이동한다. 전체 활동에 의한 나눗셈과의 관계는 총설 & \cite{NakaRushton1966,NormannPerlman1979,Carandini2012} & 측정됨 \\
7 & 센서는 변화를 전하고, 그 부호는 스파이크당 최대 정보에 맞춰져 있다(명제~\ref{prop:sensors}) & 원리는 이론적으로 제안되었다. 가장 강한 증거: 파리에서 첫 층 뉴런의 “대비 $\to$ 응답” 특성이 자연 장면 대비의 누적 분포와 일치한다. 이렇게 해서 최대 정보가 얻어진다 & \cite{Barlow1961,Laughlin1981} & 가설 \\
8 & 켜짐 센서와 꺼짐 센서는 두 전선 부호이다; 이벤트 카메라는 이를 실리콘으로 재현한다 & 실험 총설: 망막의 켜짐 채널과 꺼짐 채널은 첫 시냅스에서 이미 갈라지고 따로 차단된다. 카메라: 프레임 없는 $128\times128$픽셀, 범위 $120$\,dB, 지연 $15\,\mu\text{s}$ & \cite{Schiller1992,Lichtsteiner2008,Gallego2022} & 측정됨 \\
9 & 눈에는 빛 센서 $\sim10^8$개와 출력 뉴런 $\sim10^6$개가 있다: $\sim100$배 압축 & 사람 망막 전체 표본에서의 계수: 평균 막대세포 $92$백만 개, 원뿔세포 $4.6$백만 개; 출력(신경절) 세포는 눈에 따라 $0.7$에서 $1.5$백만 개 & \cite{Curcio1990,CurcioAllen1990} & 측정됨 \\
10 & 생쥐 눈의 출력 채널은 서른 개가 넘는다 & 생쥐: 출력 세포 $11\,000$개 이상의 이광자 칼슘 영상과 응답의 군집화, $30$개보다 뚜렷이 많은 기능 유형. 사람의 수는 측정되지 않음 & \cite{Baden2016} & 측정됨 \\
11 & 기니피그의 눈은 $\sim875\,000$ bit/s를 전송한다; 사람 눈으로 환산하면 $\sim10^7$ bit/s & 기니피그, 다전극 배열, 자연 장면 속 움직임: 세포당 $6$--$13$ bit/s, 출력 세포 $\sim10^5$개에 $\sim875\,000$ bit/s. 사람($\sim10^6$개 세포)의 $\sim10^7$ bit/s는 저자들의 환산 & \cite{Koch2006} & 측정됨 \\
12 & 귀는 거의 로그인 주파수 지도를 가진 대역 통과 필터 뱅크이다(그림~\ref{fig:filterbank}) & 기저막 최대 진동 위치는 주파수에 따라 다르다; “주파수--위치” 함수는 거의 지수적이며, 하나의 형태로 사람과 살아 있는 동물의 달팽이관을 기술한다 & \cite{Greenwood1990,Robles2001} & 측정됨 \\
13 & 공진기는 능동적이다: 센서 일부가 약한 진동을 진폭으로 수천 배($66$--$76$\,dB) 증폭하고, 소리가 커지면 증폭이 줄어든다 & 친칠라 달팽이관 기저부의 레이저 진동 측정: 약한 순음의 특성 주파수에서 등자뼈 대비 이득 $66$--$76$\,dB; 죽으면 감도가 $60$--$81$\,dB(진폭으로 $10^3$--$10^4$배) 떨어지고 압축이 사라진다. 힘의 원천은 바깥털세포 & \cite{Ruggero1997,Robles2001} & 측정됨 \\
14 & 귀의 증폭기는 자기 발진 경계 근처에 있다 & 계산: 호프 분기점 근처에서는 범위 압축, 날카로운 조율, 결합음이 불가피하다. 이것이 알려진 청각의 비선형성 전부이다. 달팽이관이 실제로 그렇게 조율되어 있다는 직접 증명은 없음 & \cite{Eguiluz2000} & 가설 \\
15 & 낮은 주파수에서 스파이크는 소리와 위상을 맞춰 나가며(기니피그에서 위상 고정은 $\sim600$\,Hz부터 약해지고 $\sim3.5$\,kHz까지 드러난다), 이로부터 방향을 찾는다; 높은 주파수에서는 위치 부호만 남는다 & 기니피그 청신경: 위상 고정은 약 $600$\,Hz에서 떨어지기 시작해 $3.5$\,kHz 위에서 사라진다; 고양이와 원숭이에서는 경계가 약 한 옥타브 높다. 두 귀의 도착 시간 차는 총설 & \cite{PalmerRussell1986,Grothe2010} & 측정됨 \\
16 & 나머지 센서(표~\ref{tab:sensors}): 후각은 수용체 유형 $\sim400$개, 뉴런당 하나, 조합 부호; 미각은 일부 ATP와 세로토닌을 통해; 촉각은 빨리, 느리게 적응하는 센서; 따뜻함과 차가움은 별도 & 생쥐: 한 수용체가 많은 냄새를, 한 냄새가 많은 수용체를 인식하며, 냄새마다 집합이 다르다. ATP 수용체가 없으면 미각 신경이 응답하지 않는다; 미뢰는 세로토닌을 방출한다. 사람 손 피부의 단일 섬유 기록: 적응 속도에 따른 네 유형. 냉각 채널은 열 채널과 다르다 & \cite{Mombaerts2004,Malnic1999,Finger2005,Huang2005,JohanssonVallbo1979,McKemy2002} & 측정됨 \\
\bottomrule
\end{longtable}}

\section{\ref{sec:recognition}장. 인식}

인식 속도, 첫 단계들의 구조, 응답의 선형 판독, 시간의 연속성에 의한 학습은 사람과 원숭이에서 측정되었다. 사슬 전체를 두 연산으로 환원하는 것과 영상에 의한 학습은 이 책의 모델로 확인했다. 착시의 설명은 근거가 탄탄한 것부터 논쟁 중인 것까지 있다.

{\footnotesize
\setlength{\tabcolsep}{3pt}\renewcommand{\arraystretch}{1.15}
\begin{longtable}{>{\raggedright}p{0.5cm}>{\raggedright}p{4.0cm}>{\raggedright}p{5.6cm}>{\raggedright}p{2.3cm}>{\raggedright\arraybackslash}p{1.7cm}}
\toprule
No. & 명제 & 실험과 측정한 것 & 출처 & 상태 \\
\midrule\endfirsthead
\toprule
No. & 명제 & 실험과 측정한 것 & 출처 & 상태 \\
\midrule\endhead
1 & 이동이 있을 때 픽셀 단위 견본 비교는 동전보다 낫지 않다; 단계 쌍~\eqref{eq:scstage}는 도형을 어디서나 인식한다 & \texttt{models/\allowbreak recognition\_models.py}, $24$개 점의 “망막”, $4000$번 시행: 뒤집힌 점의 비율 $0$, $0.1$, $0.2$에서 견본은 $0.49$, $0.51$, $0.52$; $S$와 $C$ 쌍은 $1.00$, $0.86$, $0.70$ & 이 장의 유도 & 모델 \\
2 & 그림 속 동물은 $\sim150$\,ms 만에 인식된다. 열 개쯤의 단계를 단계당 $10$--$20$\,ms씩 한 번 통과 & 사람, $20$\,ms 동안 보여 준 사진에서 “동물이 있는가” 과제: 두 답에 대한 유발 전위가 $\sim150$\,ms에 갈라진다. 원숭이: 첫 응답 시간이 단계마다 늘어난다(V1이 가장 먼저, 이어서 V2와 V4). 단계 수와 “단계당 스파이크 0개 또는 1개”는 이 수치에서 나온 추론 & \cite{Thorpe1996,Schmolesky1998} & 측정됨 \\
3 & $S$ 단계는 부분들의 AND, $C$ 단계는 위치들에 대한 최댓값이다(명제~\ref{prop:conveyor}) & 고양이, 일차 시각 피질: 단순 세포는 특정 위치의 특정 방위 가장자리에, 복합 세포는 더 넓은 영역의 같은 방위에 응답한다. 복합 세포의 세포 내 기록: 막대 두 개에 대한 응답은 많은 세포에서 두 응답 중 큰 쪽에 가깝고, 평균적으로 max 연산에 가장 가깝다. 상호 억제에 의한 최댓값은 명제~\ref{prop:wta}, 전체 활동에 의한 나눗셈은 총설 & \cite{HubelWiesel1962,Lampl2004,Carandini2012} & 측정됨 \\
4 & 사슬 위로 갈수록 조합은 커지고 복잡해지며, 응답은 위치에 점점 덜 좌우된다 & 원숭이, 자극을 “결정적 특징”까지 단순화: V4와 뒤쪽 하측두 피질 세포는 작은 수용장에서 형태, 색, 질감의 복잡한 조합에 응답한다; 앞쪽 하측두 피질에서는 수용장이 더 크다. 모델에서 $S$와 $C$의 교대가 이 그림을 재현한다 & \cite{KobatakeTanaka1994,Riesenhuber1999} & 측정됨 \\
5 & 합성곱 신경망은 이 교대를 되풀이하며 윗단계 응답을 가장 잘 예측한다; 물체는 뉴런 무리의 발화율에서 선형으로 읽힌다 & 인식을 학습했지만 뇌에 맞추지 않은 네트워크: 윗층은 원숭이 하측두 피질 뉴런의 응답을, 중간층은 V4의 응답을 예측한다. 무작위로 고른 세포 $\sim100$개의 활동에 대한 선형 분류기가 $12.5$\,ms부터의 창에서 위치와 크기가 달라도 물체와 범주를 알아본다 & \cite{Fukushima1980,Yamins2014,Hung2005} & 측정됨 \\
6 & 흔적을 가진 헵 규칙~\eqref{eq:traceinvariance}은 흔적이 $\sim20$\,ms보다 짧지 않으면 이름표 없는 영상으로 불변성을 가르친다 & 느린 특징과 흔적 규칙의 착상은 이론. \texttt{models/\allowbreak recognition\_models.py}, $C$ 뉴런 두 개, 입력 $16$개, $10$번 실행, 물체 사이 휴지 $0.3$\,s. $\tau_{\text{tr}}=5$\,ms에서 도형과의 일치는 평균 $0.52$, $10$\,ms에서 $0.56$; $20$, $50$, $200$\,ms에서 모든 실행이 $1.00$($30$\,ms에서는 열 번 중 한 번 틀림) & \cite{Foldiak1991,WiskottSejnowski2002} & 모델 \\
7 & 뇌에서 불변성은 시간의 연속성으로 학습된다 & 원숭이: 시야의 특정 위치로 눈이 도약하는 동안 물체를 몰래 바꿔치기했다; 하측두 피질 뉴런의 위치 불변성이 바꿔치기에 따라 변했고, 한 시간 만에 이미 유의했다. 이것이 시각 학습의 주된 규칙이라는 것은 가설 & \cite{LiDiCarlo2008} & 측정됨 \\
8 & 움직임: 방향 검출기, 이어서 넓은 영역에 대한 같은 AND/OR 쌍 & 토끼 망막의 방향 검출기; 원숭이 MT 영역에서 기록한 뉴런 $110$개 모두 방향 선택적이며, 움직임에 대한 응답이 V1보다 강하다 & \cite{Barlow1965,Albright1984} & 측정됨 \\
9 & 윗단계는 아래로 예측을 보내고, 위로는 오차가 올라간다 & 모델은 일차 시각 피질 수용장의 비고전적 효과를 설명한다; 개별 관찰은 부합하지만(총설) 파이프라인의 일반 구조로서는 증명되지 않음 & \cite{RaoBallard1999,Keller2018} & 가설 \\
10 & 어려운 그림은 되먹임과 여러 바퀴를 필요로 한다 & 원숭이: “어려운” 그림에서는 하측두 피질의 판정이 $\sim30$\,ms 늦게 형성된다; 이 늦은 응답은 순방향 네트워크가 잘 예측하지 못하고, 되먹임 네트워크나 더 깊은 네트워크가 더 잘 예측한다 & \cite{Kar2019} & 측정됨 \\
11 & 눈에 띄지 않는 픽셀 위조가 네트워크를 속인다; 사람의 시각은 훨씬 강하지만 무적은 아니다 & 네트워크: 특별히 고른 작은 섭동이 답을 바꾼다; “판다 $\to$ 긴팔원숭이” 예는 두 번째 논문에서. 사람에서는 짧게 보여 주면 네트워크 사이에 잘 전이되는 위조가 선택을 옮긴다 & \cite{Szegedy2014,Goodfellow2015,Elsayed2018} & 측정됨 \\
12 & 기대의 착시: 믿기 어려운 입력은 기대에 밀린다. 흐린 것은 “더 느리게” 움직이고, 드레스는 사람마다 다르게 보인다(\ref{sec:illusions}절) & 대비가 낮은 격자는 더 느리게 움직이는 것처럼 보인다. $1401$명 설문: 드레스는 세 가지 안정된 방식으로 보이고, 조명에 대한 힌트가 색을 바꾼다; 관찰자 $\sim13\,000$명에서 조명에 대한 가정이 결정한다. 느린 움직임을 기대하는 베이즈 모델이 여러 움직임 착시를 설명한다 & \cite{Thompson1982,LaferSousa2015,Wallisch2017,Weiss2002,Purves2011} & 가설 \\
13 & 이득 조절: 폭포, 기울기, 대비; 헤르만 격자는 이웃 억제가 아니다 & 토끼 망막의 방향 검출기는 오래 움직임을 보면 발화율이 떨어지고, 멈춘 뒤에는 배경 아래로 떨어진다. 기울기 착시는 이웃 활동에 의한 나눗셈 모델로 재현된다. 망막 출력 뉴런의 응답을 바꾸지 않는 격자 변형이 점을 없앤다. 이웃 억제 설명은 기각됨 & \cite{BarlowHill1963,Schwartz2009,SchillerCarvey2005} & 측정됨 \\
14 & 지연 보상: 움직이는 물체가 섬광보다 앞서 보인다 & 착시는 측정됨. 정신물리학은 움직임의 전방 연장과 지연 차이 둘 다와 모순된다; 섬광 뒤 $\sim80$\,ms의 사건에 따른 사후 설명이 제안되었다. 논쟁은 끝나지 않음 & \cite{Nijhawan1994,EaglemanSejnowski2000} & 가설 \\
15 & 정지한 “뱀”이 돈다: 지연 차이~\eqref{eq:latency}를 방향 검출기가 읽는다 & 착시는 색 없이도 일어난다; 이웃한 요소 쌍이 하나씩 착시를 만든다; 방향 선택적인 원숭이 피질 뉴런은 이 정지한 쌍에 방향성 응답을 낸다. 사람에서 회전은 미세 눈 도약과 눈 깜박임이 촉발한다 & \cite{Conway2005,OteroMillan2012} & 측정됨 \\
16 & 이중 그림: 상호 억제, 피로, 잡음; 전환은 주로 잡음이 일으킨다 & 경쟁하는 뉴런 무리 모델: 전환은 \emph{잡음}이 일으키며, 잡음이 없으면 전환도 없다; 자극 세기에 따른 전환 빈도 증가를 설명한다. 여기서 피로의 역할은 작다 & \cite{MorenoBote2007} & 가설 \\
17 & 착시의 일부는 기계가 배우는 과제의 결과이다 & 영상의 다음 프레임을 예측하도록 학습한 네트워크는 정지한 고리의 회전을 “보고” 대조 그림에서는 보지 않는다; 이미지의 잡음 제거와 선명화를 위한 네트워크는 사람처럼 색과 밝기 착시에 넘어간다 & \cite{Watanabe2018,GomezVilla2020} & 측정됨 \\
\bottomrule
\end{longtable}}

\section{\ref{sec:muscles}장. 근육으로의 출력}

출력의 구조, 곧 운동 단위, 시냅스의 안전 여유, 크기 순서대로의 동원, 근육 쌍, 신장 루프, 리듬 발생기는 직접 측정되었다. 지연 루프의 안정 조건은 정리이다. 몸의 내부 모델은 근거가 탄탄한 가설이다. 그림의 수치는 이 책의 모델에 따른 계산이다.

{\footnotesize
\setlength{\tabcolsep}{3pt}\renewcommand{\arraystretch}{1.15}
\begin{longtable}{>{\raggedright}p{0.5cm}>{\raggedright}p{4.0cm}>{\raggedright}p{5.6cm}>{\raggedright}p{2.3cm}>{\raggedright\arraybackslash}p{1.7cm}}
\toprule
No. & 명제 & 실험과 측정한 것 & 출처 & 상태 \\
\midrule\endfirsthead
\toprule
No. & 명제 & 실험과 측정한 것 & 출처 & 상태 \\
\midrule\endhead
1 & 운동 단위: \nt{ACET} 뉴런 하나가 수백에서 이천 개 남짓의 섬유 무리를 제어한다; 섬세한 조정을 하는 근육에서는 단위가 더 작다 & 사람 근육, 운동 축삭과 근육 섬유 계수: 뉴런당 섬유는 평균적으로 손의 골간근에서 약 $340$개, 완요골근에서 약 $410$개, 안쪽 장딴지근에서 약 $1930$개. 가장 작은 단위의 정확한 수는 본문에 없다. & \cite{Feinstein1955mu} & 측정됨 \\
2 & 근육 위의 시냅스는 섬유를 동작시키는 데 필요한 양의 몇 배나 되는 신경전달물질을 방출한다 & 신경근 시냅스 측정 총설: 성체 포유류의 안전 여유(방출량 대 문턱 양의 비)는 보통 $3$--$5$; 사람에서는 여유가 방출량보다 수신 측 막의 주름에 더 많이 기댄다. & \cite{WoodSlater2001} & 측정됨 \\
3 & 연축~\eqref{eq:twitch}은 $50\ms$ 정도의 시간 $T_c$ 동안 커진다 & 수의적 힘을 낼 때 사람 손의 단일 운동 단위, 단위의 스파이크로 평균한 힘: 상승 시간 $30$--$100\ms$, 단위의 $80\,\%$가 $70\ms$ 미만. 함수 $(t/T_c)e^{1-t/T_c}$는 단위 풀 모델의 근사. & \cite{MilnerBrown1973rec, Fuglevand1993} & 측정됨 \\
4 & 연축의 합~\eqref{eq:forcesum}: $5$, $15$, $40$ Hz에서 힘 $0.2$--$1.1F_1$, $1.8$--$2.2F_1$, 약 $5.4F_1$(그림~\ref{fig:tetanus}) & \texttt{models/\allowbreak muscle\_\allowbreak models.py}에 따른 계산, $T_c=50\ms$: 마지막 $0.3\,\text{s}$에서 $0.21$--$1.10$, $1.76$--$2.19$, $5.28$--$5.49\,F_1$. 선형 합은 단순화: 실제 근육의 포화는 모델에 없다. & \texttt{models/\allowbreak muscle\_\allowbreak models.py} & 모델 \\
5 & 단위는 크기 순서대로 켜진다; 가장 약한 단위는 가장 강한 단위보다 백 배 약하다 & 고양이 배쪽 뿌리: 반사 입력이 커질 때 뿌리에서 스파이크가 작은 뉴런, 곧 작은 뉴런이 먼저 방전한다. 사람 손의 골간근: 단위의 연축 힘은 $0.1$에서 $10$~g이고, 단위가 켜지는 힘에 거의 선형으로 커진다. & \cite{Henneman1957, MilnerBrown1973rec} & 측정됨 \\
6 & 힘의 단계는 힘에 비례한다, $\Delta F/F\approx q-1$~\eqref{eq:sizeprinciple}: 부동소수점 & 힘의 등비수열에서 본문이 유도. 실험과 부합: 큰 힘에서는 같은 힘 증가에 새로 켜지는 단위가 훨씬 적다. & 이 장의 유도; \cite{MilnerBrown1973rec} & 이론 \\
7 & 힘의 퍼짐은 힘 자체에 비례해 커진다 & 엄지 근육의 수의적 등척성 힘: 힘의 표준편차는 평균 힘에 따라 선형으로 커진다; 전기 자극으로 같은 힘을 낼 때는 퍼짐이 일정하다. 단위 풀 모델: 힘 순서대로의 동원이 선형 증가를 준다. & \cite{Jones2002sdn} & 측정됨 \\
8 & 움직임의 매끄러움은 신호 의존 잡음의 결과이다 & 모델: 이런 잡음에서 끝점 퍼짐을 최소화하는 궤적이 측정된 눈과 팔의 궤적을 재현한다. 이 원인을 다른 원인과 구별하는 직접 실험은 없다. & \cite{HarrisWolpert1998} & 가설 \\
9 & 근육 쌍의 힘 차가 회전력을, 합이 강성을 정한다~\eqref{eq:antagonists} & 공동 긴장에 의한 임피던스 제어 이론. 사람 팔: 작은 밀기로 잰 각 관절의 강성은 그 관절의 회전력에 대략 비례한다; 공동 긴장의 비율이 다르면 손 강성 타원의 모양과 방향이 바뀐다. & \cite{Hogan1984, GomiOsu1998stiff} & 측정됨 \\
10 & 신장 센서는 자기 \nt{ACET} 뉴런을 흥분시키고, 억제성 중개 뉴런을 거쳐 길항근을 끈다(그림~\ref{fig:antagonists}) & 고양이 척수: 신장 센서 섬유로 흥분되는 단일 중개 뉴런이 길항근 \nt{ACET} 뉴런에 단일 시냅스 억제 전위를 만든다, 시냅스 지연 $0.28$--$0.42\ms$. 중개 뉴런의 신경전달물질(글리신)은 이 실험에서 확인하지 않았다. & \cite{Jankowska1972Ia} & 측정됨 \\
11 & 루프 지연: 척수 루프 $25$--$30\ms$, 피질을 거치면 1.5--3배, 눈을 거치면 $100$--$200\ms$ & 사람 팔, 관절 밀기: 근육의 짧은 응답은 $20$--$45\ms$, 긴 응답은 $45$--$105\ms$, 수의적 응답은 $120$--$180\ms$ 뒤. 움직이는 동안 보이는 손가락 위치를 옮기면: 보정은 평균 $160\ms$ 뒤. & \cite{Pruszynski2008rapid, SaundersKnill2003vis} & 측정됨 \\
12 & $\dd x/\dd t=-Gx(t-d)$는 $Gd<\pi/2$일 때 안정하다~\eqref{eq:delaystable}; 경계에서 주파수는 $1/(4d)$ & 지연 방정식의 근에 관한 정리. \texttt{models/\allowbreak muscle\_\allowbreak models.py}로 확인, $d=30\ms$: $3\,\text{s}$까지의 진폭은 초당 $G=40$에서 $3\cdot10^{-6}$, $G=50$에서 $0.13$, $G=55$에서 $38$(경계 $52$). & \cite{Hayes1950} & 이론 \\
13 & 생리적 떨림 $8$--$12$~Hz; 신장 루프는 그 원천 가운데 하나 & 측정 총설: 약 $10$~Hz 대역의 미세 떨림은 건강한 사람에게 있다; 그 기원은 복합적이다. 중추 진동, 단위 방전의 성질, 기계적 공진, 반사 루프의 공진이 있고, 조건에 따라 우세한 것이 다르다. & \cite{McAuleyMarsden2000} & 가설 \\
14 & 뇌는 몸의 내부 모델로 명령의 결과를 예측한다 & 사람이 물체를 손가락으로 집어 들고 팔을 움직인다: 관성, 점성, 탄성 부하에서 쥐는 힘이 루프 지연 없이 부하와 동시에 변한다. 곧 부하가 예측된다. 총설이 이런 실험을 정리한다; 뉴런에서 모델 자체를 직접 관찰한 것은 없다. & \cite{FlanaganWing1997grip, WolpertGhahramani2000} & 가설 \\
15 & 걷기, 숨쉬기, 씹기의 리듬은 일정한 입력에서 국소 네트워크가 만든다 & 실험 총설: 분리된 신경계(척수, 신경절)는 리듬적 입력도 근육의 되먹임도 없이 리듬 운동 패턴을 낸다. & \cite{MarderBucher2001} & 측정됨 \\
16 & 피로를 동반한 상호 억제~\eqref{eq:halfcentre}는 주기 $0.84\,\text{s}\approx2\tau_a$의 리듬을 준다(그림~\ref{fig:halfcentre}) & 정리: 상호 억제와 적응을 가지며 일정한 입력에서 안정 상태가 없는 네트워크는 반드시 진동한다. \texttt{models/\allowbreak muscle\_\allowbreak models.py}에 따른 계산($I=1$, $\beta=2$, $b=2$, $\tau=20\ms$, $\tau_a=0.4\,\text{s}$): 주기 $0.84\,\text{s}$. 실제 발생기가 바로 이렇게 만들어져 있다는 것은 보이지 않았다. & \cite{Matsuoka1985osc} & 모델 \\
\bottomrule
\end{longtable}}

\section{\ref{sec:pain}장. 통증}

통증 센서, 두 전선, 척수의 관문, 하행 음량 노브, 감작, 주의 포획은 직접 측정되었다. 관문과 감작의 식은 이 책의 단순화된 모델이다. 기계가 경보 음량을 고르는 규칙은 가설이다.

{\footnotesize
\setlength{\tabcolsep}{3pt}\renewcommand{\arraystretch}{1.15}
\begin{longtable}{>{\raggedright}p{0.5cm}>{\raggedright}p{4.0cm}>{\raggedright}p{5.6cm}>{\raggedright}p{2.3cm}>{\raggedright\arraybackslash}p{1.7cm}}
\toprule
No. & 명제 & 실험과 측정한 것 & 출처 & 상태 \\
\midrule\endfirsthead
\toprule
No. & 명제 & 실험과 측정한 것 & 출처 & 상태 \\
\midrule\endhead
1 & 높은 문턱값: 통증 센서마다 가열 문턱값은 $41^\circ$에서 $46$--$48^\circ$까지 & 단일 섬유 기록. 원숭이 피부: 느린(C) 센서의 중앙 문턱값 $41^\circ$, 빠른 제2형 센서 $46^\circ$, 제1형은 $53^\circ$ 이상. 사람 피부: 압력에도 응답하는 C 센서 $40.7^\circ$, 압력에 응답하지 않는 것 $48^\circ$. & \cite{Treede1995heat, Weidner1999cfib} & 측정됨 \\
2 & 통증 센서는 촉각 센서보다 훨씬 약하게 적응하고, 가벼운 손상 뒤에는 더 민감해진다; 강한 가열 뒤의 약화는 한 유형에서 측정됨 & 피부의 가벼운 화상($50^\circ$, $100\,\text{s}$): $5$--$10$~분 뒤 사람의 통증 문턱값과 원숭이 C 센서의 문턱값이 $2$--$6^\circ$ 내려간다; 다른 피부 센서에는 이런 이동이 없다. 빠른 제2형 센서는 강한 가열 뒤 응답이 억제된다. 촉각 센서와의 적응 비교는 일반적인 추정으로, 여기서 별도 실험으로 뒷받침되지는 않는다. & \cite{LaMotte1982hyper, Treede1995heat} & 측정됨 \\
3 & 두 전선: 빠른 전선 $5$--$30$~m/s, 느린 전선 $0.5$--$2$~m/s; 도착 시간 $t=L/v$~\eqref{eq:paintime} & 전도 속도: 원숭이의 빠른 통증 센서는 평균 $15$와 $25$~m/s(두 유형); 사람의 C 센서는 $0.8$--$1.0$~m/s. $L=1$~m이면 수십 밀리초와 약 1초. & \cite{Treede1995heat, Weidner1999cfib} & 측정됨 \\
4 & 통증은 두 번 온다: 먼저 빠르고 정확하게, 다음에 둔하고 오래 & 사람 아래팔 가열에 대한 반응 시간: 온도가 오르면 $1100$에서 $400\ms$로; 털 있는 피부의 강한 가열은 이중 통증을 주고, 손바닥은 주지 않는다. 첫 통증은 $6$~m/s보다 빠른 섬유만 나를 수 있다. 잠복기로 보면 손바닥에 없는 빠른 제2형 센서가 이에 해당한다. 역할 분담(“어디”와 “얼마나 나쁜지”)은 해석. & \cite{CampbellLaMotte1983first, Treede1995heat} & 측정됨 \\
5 & 관문: 척수 출력 뉴런이 통증과 촉각을 받고, 억제성 중개 뉴런은 촉각에 흥분한다(그림~\ref{fig:gate}) & 관문 도식은 이론으로 제안되었다. 생쥐, 유전적 표지와 뉴런 집단 제거: 소마토스타틴을 가진 흥분성 뉴런은 기계적 통증에 필요하다; 다이노르핀을 가진 억제성 뉴런은 촉각 센서의 입력이 그 뉴런에 닿지 못하게 한다. 이 뉴런이 없으면 가벼운 접촉이 통증을 일으킨다. & \cite{MelzackWall1965, Duan2014} & 측정됨 \\
6 & 촉각은 통증을 누그러뜨린다 & 사람, 손의 레이저 통증 펄스: 동시에 접촉하면 통증 평가와 펄스 에너지 구별 능력이 낮아진다; 한 피부 분절 안에서 접촉 지점과 통증 지점 사이 거리가 멀수록 효과가 약해진다. & \cite{Mancini2014touch} & 측정됨 \\
7 & 관문 이론의 가정: 강한 통증이 스스로 억제성 중개 뉴런을 꺼서 관문이 활짝 열린다 & 본문에서 원래 이론의 가정으로 표시했다. 생쥐 연구는 관문이 촉각을 통증 채널로 들이지 않는 방식을 기술한다; 통증에 의한 중개 뉴런의 꺼짐은 거기서 보이지 않았다. & \cite{MelzackWall1965} & 가설 \\
8 & 관문~\eqref{eq:gate}: 촉각 없이 문턱값 $0.18$, 촉각이 있으면 $0.86$, $g=2$에서 $0.11$; $\nu=1$에서 촉각은 출력을 $1$에서 $0.25$로 낮추고, $\nu=2$에서는 효과가 없다; 촉각 자체는 통과하지 못한다(그림~\ref{fig:gatecurves}) & 본문의 가중치로 \texttt{models/\allowbreak pain\_\allowbreak gate.py}에 따라 계산하면 이 수치가 모두 재현된다. & \texttt{models/\allowbreak pain\_\allowbreak gate.py} & 모델 \\
9 & 뇌간에서 내려오는 하행 경로는 관문의 이득을 양쪽으로 바꾼다 & 쥐, 연수, 꼬리 가열 중 기록: 어떤 뉴런은 꼬리를 빼기 전에 방전하고(“켜짐”), 다른 뉴런은 조용해진다(“꺼짐”); 이 영역의 약한 자극은 반사를 억제한다. 총설: 같은 회로가 통증 전달을 촉진하기도 억제하기도 하며, 오피오이드는 이를 통해 작용한다. & \cite{Fields1983rvm, Fields2004} & 측정됨 \\
10 & 통증이 줄어들 것이라는 기대는 오피오이드 채널을 통해 통증을 누그러뜨린다 & 급성 통증이 있는 사람들: 기대로 통증이 줄어든 사람에서는 오피오이드 수용체 차단이 통증을 다른 사람들 수준으로 되돌렸다; 다른 사람들에서는 차단이 통증을 바꾸지 않았다. & \cite{Levine1978} & 측정됨 \\
11 & 뇌는 중단의 이득에 따라 경보 음량을 고른다 & 음량 선택 규칙을 측정한 실험은 없다. & --- & 가설 \\
12 & 감작: 손상 뒤 관문 출력이 커지고 문턱값이 떨어지며, 일부는 척수에서 일어난다; 손상 자극이 멈춘 뒤에도 유지된다 & 쥐: 피부 손상 뒤 굽힘 반사의 문턱값이 떨어지고 응답이 커진다; 전기생리학적 분석은 증가의 일부가 척수 자체에서 생김을 보여 준다. 손상 자극은 자극 자체보다 오래가는 감도의 장기 변화를 일으킨다. 정확한 회복 시간은 본문에 없다. & \cite{Woolf1983} & 측정됨 \\
13 & 감작~\eqref{eq:sensitization}은 양의 되먹임이다; 루프가 강하면 경보는 치유 뒤에도 걸린 채 남을 수 있다 & 모델~\eqref{eq:sensitization}에서 나온다(장 끝의 연습문제). 치유 뒤 지속되는 통증이 바로 이런 형태로 걸린 루프임을 보인 실험은 없다. & 장 안의 유도 & 가설 \\
14 & 통증은 음의 보상이며, 그것으로 하는 학습은 시간차 오차를 따른다 & 영상 장치, 사람, 통증 전조의 사슬: 배쪽 선조체와 앞쪽 섬엽 피질의 활동이 시간차 학습 신호와 일치한다. 원숭이: \nt{DOPA} 뉴런 일부는 불쾌한 공기 분사와 그 전조에 흥분하고, 다른 일부는 억제된다. & \cite{Seymour2004, Matsumoto2009} & 측정됨 \\
15 & 통증은 진행 중인 작업을 중단시킨다 & 사람, 전기 통증 자극과 소리 시험: 통증 시작 $100\ms$ 뒤의 시험 응답은 뚜렷이 느려지고, $1.5\,\text{s}$ 뒤에는 대조군과 차이가 없다. 총설은 이런 실험을 주의 포획 모델로 정리한다. & \cite{Crombez1996disrupt, EcclestonCrombez1999} & 측정됨 \\
16 & 회피 반사는 척수에서 닫힌다 & 쥐: 후각 깊은 층의 뉴런이 개별 근육의 회피 반사 공간 패턴을 되풀이한다; 경수에서 역방향으로 흥분시킬 수 없으므로 척수 중개 뉴런이다. & \cite{Schouenborg1995wr} & 측정됨 \\
\bottomrule
\end{longtable}}

\section{\ref{sec:motorlearning}장. 운동 학습}

최소 제곱 규칙과 개별 오차 전선의 이점은 정리이다. 오차 전선 회로의 구조, 겹칠 때의 약화, 지연에 맞춘 적격 흔적, 후효과와 기술의 자발적 복귀는 측정되었다. 회로 전체가 지도 학습으로 동작하며 내부 모델을 만든다는 것은 유력한 가설이다. 두 과정 모델의 수치는 이 책의 모델에 따른 계산이다.

{\footnotesize
\setlength{\tabcolsep}{3pt}\renewcommand{\arraystretch}{1.15}
\begin{longtable}{>{\raggedright}p{0.5cm}>{\raggedright}p{4.0cm}>{\raggedright}p{5.6cm}>{\raggedright}p{2.3cm}>{\raggedright\arraybackslash}p{1.7cm}}
\toprule
No. & 명제 & 실험과 측정한 것 & 출처 & 상태 \\
\midrule\endfirsthead
\toprule
No. & 명제 & 실험과 측정한 것 & 출처 & 상태 \\
\midrule\endhead
1 & 규칙~\eqref{eq:lms}는 평균적으로 오차 제곱의 기울기를 따라 내려간다 & 적응 필터 이론의 결과: 입력과 출력 오차에 비례하는 보정은 평균 제곱 오차의 확률적 경사 하강이다. & \cite{WidrowHoff1960} & 이론 \\
2 & 출력마다 오차 전선이 있으면 속도는 출력 수와 무관하다; 하나의 공통 수로는 잡음 내는 뉴런 수에 따라 떨어진다(명제~\ref{prop:teachers}) & 선형 네트워크의 학습 곡선: 뉴런의 무작위 섭동에 의한 학습은 정확한 기울기보다 섭동받는 뉴런 수만큼 느리다. & \cite{Werfel2005} & 이론 \\
3 & 오차 전선 회로는 지도 학습으로 동작한다(명제~\ref{prop:errorwire}) & 회로 구조에서 유도한 이론. 7--10행의 실험은 이와 부합한다; 회로 전체가 바로 이렇게 동작한다는 것은 증명되지 않았다. & \cite{Marr1969, Albus1971} & 가설 \\
4 & 확장층: 작은 \nt{GLUT} 뉴런 약 $7\cdot10^{10}$개, 뇌의 나머지 전부보다 많다 & 녹인 조직 현탁액에서 핵 계수: 사람 소뇌에 뉴런 $69\cdot10^9$개, 뇌 전체는 $86\cdot10^9$개. 입체해석학: 노년 남성의 소뇌에 작은 뉴런 $101\cdot10^9$개. & \cite{Azevedo2009, Andersen1992cereb} & 측정됨 \\
5 & 입력의 차원을 높이면 원하는 관계가 선형이 된다 & 분할 수에 관한 정리: 문턱값을 가진 입력 $n$개의 가중합은 일반 위치에 있는 벡터 약 $2n$개의 임의의 표지를 구현한다. 소뇌가 바로 이것을 쓴다는 것은 3행 가설의 일부이다. & \cite{Cover1965} & 이론 \\
6 & 출력 \nt{GABA} 뉴런 약 $3\cdot10^7$개, 각각 $10^5$개 정도의 입력 & 사람 소뇌의 입체해석학: 출력 뉴런 $30.5\cdot10^6$개. 전자 현미경, 쥐: 출력 뉴런 하나에 확장층 입력 약 $1.75\cdot10^5$개. 출력 뉴런의 억제성 신경전달물질은 총설. & \cite{Andersen1992cereb, NapperHarvey1988, Ito2001} & 측정됨 \\
7 & 출력 뉴런마다 정확히 하나의 오차 전선, 초당 약 한 번, 매번 강력한 폭발 & 기록 총설: 각 출력 뉴런은 등반하는 \nt{GLUT} 입력 하나를 받으며, 이것이 $1$~Hz 정도의 빈도로 복합 스파이크를 일으킨다. & \cite{Ito2001} & 측정됨 \\
8 & 폭발의 확률은 움직임 오차의 방향에 따라 다르다 & 원숭이, 소뇌의 안구 운동 부위, 도약 안구 운동 적응: 시각 오차의 방향이 복합 스파이크의 확률을, 크기가 그 시각을 정한다; 폭발은 다음 시도의 단순 스파이크를 바꾸고 움직임을 세포의 선호 오차 방향으로 옮긴다. & \cite{Herzfeld2018} & 측정됨 \\
9 & 오차 폭발이 활성 입력과 겹치면 그 입력이 약해진다($-\eta e_i x_j$) & 실험 총설: 확장층 입력과 오차 전선을 함께 자극하면 그 입력이 장기적으로 약해진다; 한 입력만 자극하면 그렇지 않다. & \cite{Ito2001} & 측정됨 \\
10 & 흔적의 길이는 오차 지연에 맞춰져 있다: 지연이 $100\ms$이면 $100\ms$ 앞서 활성이던 입력이 가장 강하게 약해진다 & 절편과 살아 있는 생쥐: 안구 운동 학습을 맡는 부위에서 약화는 입력과 오차 사이 간격 약 $120\ms$에 좁게 맞춰져 있다. 이 과제에서 오차 신호가 오는 데 걸리는 시간이다; 다른 부위에서는 세포마다 다른 간격에 맞춰져 있다. & \cite{Suvrathan2016} & 측정됨 \\
11 & 회로는 몸의 내부 모델을 학습하는 적응 필터~\eqref{eq:adaptivefilter}이다 & 회로 구조에서 유도한 모델. 학습된 필터를 몸의 전달 함수로 측정한 직접 실험은 없다. & \cite{Fujita1982} & 가설 \\
12 & 예측은 자기 움직임의 결과를 빼므로 스스로를 간지럽힐 수 없다 & 사람: 자기 접촉은 같은 남의 접촉보다 덜 간지럽게 느껴진다; 영상 장치에서 체감각 피질은 자기 접촉에 더 약하게 응답하고, 움직임이 피부에 닿을 때 소뇌는 덜 활성이다. 예측을 빼는 설명은 가설. & \cite{Blakemore1998} & 가설 \\
13 & 외부 힘이 있으면 빗나감이 줄고, 힘을 없애면 팔이 반대쪽으로 빗나간다 & 사람이 힘의 장 속에서 로봇 손잡이를 잡고 뻗는다: 궤적이 직선으로 돌아온다; 장을 없애면 후효과는 처음 왜곡의 거의 거울상이다; 훈련하지 않은 공간 영역으로도 전이된다. & \cite{ShadmehrMussaIvaldi1994} & 측정됨 \\
14 & 두 과정~\eqref{eq:tworate}이 학습한다; 반대 섭동 뒤에 기술이 스스로 돌아온다 & 사람, 힘의 장, 이어서 짧은 반대 장과 오차를 고정한 시도: 보정이 스스로 옛 값으로 돌아간다; 두 과정 모델은 이를 재현하고 한 과정 모델은 재현하지 못한다. & \cite{Smith2006} & 측정됨 \\
15 & 그림~\ref{fig:tworate}의 수치: $200$번 움직임에 $0.84$(느린 과정 $0.63$); 반대 $6$번; 빠른 과정 $-0.53$, 느린 과정 $0.46$; $40$번 움직임에 $0.37$까지 복귀 & \texttt{models/\allowbreak motor\_\allowbreak learning.py}에 따른 계산, $A_f=0.92$, $B_f=0.1$, $A_s=0.996$, $B_s=0.02$: $0.840$($0.634$와 $0.205$), 움직임 $6$번, $-0.531$과 $0.457$, $40$번째 움직임에서 최고 $0.371$. & \texttt{models/\allowbreak motor\_\allowbreak learning.py} & 모델 \\
16 & 몸이 여러 속도로 변하므로 두 과정이 필요하다 & 이론: 수명이 여럿인 교란에 대한 최적 추정은 시간 상수가 다른 여러 필터를 주며, 자발적 복귀와 빨라진 재학습을 설명한다. & \cite{Kording2007} & 가설 \\
\bottomrule
\end{longtable}}

\section{\ref{sec:chemoutput}장. 화학적 출력}

심장으로 가는 두 가닥의 전선과 그 속도 차이, 호르몬 캐스케이드의 음의 되먹임, 분량 빈도에 의한 부호, 일주기 발진기와 빛에 의한 맞춤은 직접 측정되었다. $K<8$이라는 경계는 본문에서 유도한 제어 이론의 정리다. 캐스케이드의 되먹임 자체가 맥동을 만든다는 것은 이를 강하게 뒷받침하는 실험이 있는 가설이다.

{\footnotesize
\setlength{\tabcolsep}{3pt}\renewcommand{\arraystretch}{1.15}
\begin{longtable}{>{\raggedright}p{0.5cm}>{\raggedright}p{4.0cm}>{\raggedright}p{5.6cm}>{\raggedright}p{2.3cm}>{\raggedright\arraybackslash}p{1.7cm}}
\toprule
No. & 명제 & 실험과 측정한 것 & 출처 & 상태 \\
\midrule\endfirsthead
\toprule
No. & 명제 & 실험과 측정한 것 & 출처 & 상태 \\
\midrule\endhead
1 & 심장 페이스메이커는 스스로 1분에 약 $100$번 뛴다. 안정 시에는 브레이크가 밟혀 있어 박동수가 보통 $60$--$70$ 정도다 & 사람, 심장으로 가는 두 전선을 모두 차단: 고유 박동수는 안정 시 박동수보다 높고 나이가 들수록 줄며, 성인에서 분당 약 $100$회다. 따라서 안정 시에는 브레이크가 우세하다. 안정 시 박동수 $60$--$70$은 전형적인 값으로 따로 인용하지 않았다. & \cite{JoseCollison1970ihr} & 측정됨 \\
2 & 가속 페달 \nt{NORA}와 브레이크 \nt{ACET}은 저역 통과 필터이며 브레이크가 더 빠르다~\eqref{eq:gasbrake} & 개, 심장의 미주 신경 또는 교감 신경을 주파수 변조 자극하고 심방 박동수까지의 전달 함수를 측정: 두 경로 모두 저역 통과 필터이고, 교감 경로는 차단 주파수가 훨씬 낮으며 약 $1.7\,\text{s}$의 순수 지연이 더해진다. 특성은 평균 긴장도에 따라 달라지므로 선형식~\eqref{eq:gasbrake}은 근사다. & \cite{Berger1989} & 측정됨 \\
3 & 브레이크가 빠른 것은 \nt{ACET}이 2차 전달자 없이 칼륨 채널을 열기 때문이고, \nt{NORA}는 연쇄 반응을 거친다 & 심방 세포의 막 조각: 아세틸콜린이 여는 칼륨 채널은 수용체와 결합한 G 단백질의 $\beta\gamma$ 소단위가 세포 안 2차 전달자 없이 연다. cAMP를 거치는 \nt{NORA} 경로는 여기서 인용하지 않았다. & \cite{Logothetis1987gbg} & 측정됨 \\
4 & 혈액은 약 1분 만에 몸을 한 바퀴 돈다. 혈류 속 \nt{ADRE}는 수용체가 있는 모든 기관에 닿는다 & 일반적인 크기 수준의 추정. 본문에서도 그렇게 서술했고 출처는 인용하지 않았다. & --- & 추정 \\
5 & 호르몬 캐스케이드는 음의 되먹임으로 감싸여 있다. 최종 호르몬이 앞 단계를 억제한다 & 실험 리뷰: 부신 피질 호르몬은 뇌하수체의 ACTH 분비와 시상하부의 방출 호르몬 분비를 빠른, 중간, 느린 세 시간 범위에서 억제한다. & \cite{KellerWood1984fb} & 측정됨 \\
6 & 같은 단 세 개는 $K<8$에서 안정하고~\eqref{eq:cascadestable}, 경계에서의 주기는 $2\pi\tau/\sqrt3\approx3.6\tau$ & 본문에서 유도: 주파수 $\sqrt3/\tau$에서 세 단은 위상을 $180^\circ$ 늦추고 진폭을 $8$분의 1로 줄인다. & 본문에서 유도 & 이론 \\
7 & 농도 오차는 $1/(1+K)$이므로 이런 조절기는 약 $10\,\%$보다 정확하게 유지하지 못한다(명제~\ref{prop:cascade}) & 비례 되먹임에 대한 6행의 따름정리. 실제 축에는 여러 단계와 여러 시간 범위에 되먹임이 있으므로(5행), 이 경계는 모델~\eqref{eq:cascade}에 대한 것이며 측정된 것이 아니다. & 본문에서 유도 & 이론 \\
8 & $K=4$와 $7$에서 농도가 자리 잡고, $K=12$에서는 주기 $3.7$--$3.9\,\tau$의 일정한 맥동(그림~\ref{fig:cascade}) & \texttt{models/\allowbreak chem\_\allowbreak models.py}로 계산: $K=12$에서 연속된 주기 $3.9$, $3.7$, $3.8$, $3.7\,\tau$; $K=4$에서 계산 끝의 진폭 $0.003$. & \texttt{models/\allowbreak chem\_\allowbreak models.py} & 모델 \\
9 & 스트레스 호르몬은 약 1시간에 한 번 맥동으로 분비된다. 맥동은 별도의 발진기 없이 되먹임 자체가 만든다 & 쥐, 방출 호르몬을 일정하게 주입: ACTH와 코르티코스테론이 생리적인, 거의 1시간 주기로 오르내린다. 맥동에 시상하부의 리듬 입력은 필요 없다. 지연된 되먹임이 있는 뇌하수체--부신 모델이 이런 진동을 낸다. & \cite{Walker2012glc, Walker2010} & 가설 \\
10 & 수신자는 농도가 아니라 분량의 빈도를 읽는다 & 생식샘 자극 호르몬 방출 호르몬을 스스로 분비하지 못하는 원숭이: 일정한 주입은 뇌하수체 호르몬 분비를 회복시키지 못하고, 1시간에 한 번 분량으로 주면 회복된다. 다시 일정한 주입으로 바꾸면 반응이 다시 사라진다. 고역 통과 필터로서의 수용체 피로는 해석이다. & \cite{Belchetz1978} & 측정됨 \\
11 & 분량의 빈도가 다르면 같은 샘의 서로 다른 세포에 다르게 작용한다 & 같은 원숭이: 분량을 1시간에 $2$--$5$회로 늘리면 두 뇌하수체 호르몬이 모두 줄고, $3$시간에 1회로 줄이면 하나(LH)는 줄고 다른 하나(FSH)는 늘어난다. 따라서 둘의 비는 빈도가 정한다. & \cite{Wildt1981gnrh} & 측정됨 \\
12 & 일주기 리듬은 몇 시간 지연되는 세포 안의 음의 되먹임이 만든다 & 리뷰: 시계 유전자의 단백질이 지연을 두고 자기 유전자의 전사를 억제한다. 전사와 번역으로 이루어진 루프가 약 하루의 주기를 낸다. & \cite{TakahashiJS2017} & 측정됨 \\
13 & 주 발진기는 수만 개의 뉴런으로 이루어진 무리이며 몸의 나머지를 동기화한다 & 리뷰: 시교차상핵이 주 일주기 발진기다. 배양한 개별 뉴런은 독립된 발진기이고, 네트워크가 이들을 동기화한다. 생쥐에서는 한쪽에 약 $10^4$개의 뉴런이 있다. & \cite{Welsh2010scn} & 측정됨 \\
14 & 사람의 고유 주기는 $24.18$시간 & 하루와 분리된 일정 아래 어두운 빛에서 지낸 사람: 멜라토닌, 체온, 코르티솔 리듬의 주기는 젊은 사람과 노인 모두 평균 $24.18$시간이고 편차가 좁다. & \cite{Czeisler1999} & 측정됨 \\
15 & 빛은 위상 고정 루프처럼 발진기를 맞춘다~\eqref{eq:pll}. 하루 약 $11$~min의 이동이 필요하다 & 사람, 하루 중 여러 시각에 $6.7$시간의 밝은 빛을 한 번 비춤: 진폭 $5$시간의 1형 위상 반응 곡선. 체온 최저점 이전에는 지연, 이후에는 앞당김. 식~\eqref{eq:pll}은 표준 모델이고, $11$~min $=0.18$시간은 산술이다. & \cite{Khalsa2003prc} & 측정됨 \\
16 & 빛이 없으면 맞춤이 없고 리듬은 고유 주기로 흘러간다 & 빛을 전혀 지각하지 못하고 일상 사회에서 사는 전맹인: 약 절반에서 멜라토닌과 코르티솔 리듬이 하루와 맞지 않는 자기 주기로 간다. & \cite{Sack1992blind} & 측정됨 \\
\bottomrule
\end{longtable}}

\section{\ref{sec:debugging}장. 기계 디버깅}

전극이 무엇을 듣는지, 활동의 저차원성, 딱딱한 배열을 쓴 인터페이스의 성능, 뇌와 디코더의 공동 학습, 자극으로 생기는 시각의 점은 동료 심사를 거친 실험에서 측정되었다. 데이터 흐름의 추정은 본문의 산술이다. 사람에게 이식한 Neuralink 장치에 대해서는 확인한 범위에서 주로 회사 자체가 데이터를 발표했다.

{\footnotesize
\setlength{\tabcolsep}{3pt}\renewcommand{\arraystretch}{1.15}
\begin{longtable}{>{\raggedright}p{0.5cm}>{\raggedright}p{4.0cm}>{\raggedright}p{5.6cm}>{\raggedright}p{2.3cm}>{\raggedright\arraybackslash}p{1.7cm}}
\toprule
No. & 명제 & 실험과 측정한 것 & 출처 & 상태 \\
\midrule\endfirsthead
\toprule
No. & 명제 & 실험과 측정한 것 & 출처 & 상태 \\
\midrule\endhead
1 & 전극은 반경 약 100마이크로미터 안의 뉴런, 보통 한 개에서 몇 개의 스파이크를 듣고, 모양으로 구별한다 & 쥐, CA1 영역, 세포 안과 세포 밖 동시 기록: 세포 밖 스파이크의 모양은 세포 안 스파이크의 폭과 진폭을 따른다. 추정: 테트로드는 뉴런 $60$--$100$개를 구별할 수 있지만 실제로는 약 여섯 개를 골라낸다. & \cite{Henze2000extra} & 측정됨 \\
2 & 뇌에는 뉴런 $8.6\cdot10^{10}$개가 있고, 프로브는 대략 1억분의 1을 듣는다 & 녹인 조직의 현탁액에서 핵을 셈: 사람 뇌의 뉴런 $86\cdot10^9$개. 비율은 본문의 산술. & \cite{Azevedo2009} & 측정됨 \\
3 & 운동 피질 뉴런 수백 개의 활동은 공통 변수 열 개 남짓에 들어간다 & 원숭이 운동 피질 기록의 리뷰: 집단 활동은 많은 뉴런이 공유하는 소수의 잠재 변수로 기술된다. & \cite{Gallego2017} & 측정됨 \\
4 & 이웃의 잡음은 상관되어 있고, 뉴런 100개는 1000개와 거의 같은 양을 전달한다(명제~\ref{prop:pooling}) & 원숭이 시각 피질: 이웃 뉴런의 잡음 상관 계수는 약 $0.1$이며, 평균의 이득은 뉴런 100개에서 포화된다. 정보를 제한하는 상관의 이론. & \cite{Zohary1994, MorenoBote2014} & 측정됨 \\
5 & 원시 흐름 $2\cdot10^8$~bit/s 대 스파이크 $8\cdot10^4$~bit/s~\eqref{eq:probedata} & 스파이크 흐름의 용량~\eqref{eq:spikeentropy}에 따른 본문의 산술. 디지털화 매개변수는 실제 값이다: Neuralink 칩은 채널당 $19.3$~kHz, $10$비트. & 이 장의 유도; \cite{Rieke1997, Musk2019} & 이론 \\
6 & Neuralink의 첫 시스템: 칩이 증폭하고 디지털화하며, 원시 흐름은 케이블로 나가고, 스파이크는 외부 프로그램이 골라낸다. 칩 위의 스파이크 검출, 밀폐된 본체, 무선 통신과 무선 전력은 회사 발표에 따르면 사람용 장치에 있다 & 회사 논문: 원시 흐름 전체가 USB-C 케이블로 나가고, 스파이크는 칩 밖의 프로그램이 실시간으로 골라낸다. 탑재 압축, 무선 전력과 전송은 임상 장치의 계획으로 언급되었다. 사람에게 이식한 장치는 회사 발표뿐이다. 칩 위에서 스파이크를 골라내야 한다는 것은 \eqref{eq:probedata}에서 나온 본문의 결론이다. & \cite{Musk2019}; Neuralink 보고, 동료 심사 논문 없음 & 측정됨(회사 논문); 사람은 회사 보고 \\
7 & 실 $96$가닥에 전극 최대 $3072$개, 실마다 $32$개; 실은 로봇이 혈관을 피해 꽂는다 & 회사 논문: 실 $96$가닥에 전극 $3072$개, 실마다 $32$개; 로봇이 1분에 실 $6$가닥(전극 $192$개)을 꽂는다; 배열을 오래 이식한 쥐에서 전극의 최대 $70\,\%$가 스파이크를 듣는다. & \cite{Musk2019} & 측정됨(회사 논문) \\
8 & 사람에서는 실 $64$가닥에 약 1000채널; 실 일부가 밀려났다; 커서는 약 $10$~bit/s & 회사 발표뿐. PubMed 검색에서 사람 대상 결과를 담은 동료 심사 논문을 찾지 못했다. 본문의 단서와 일치한다. & Neuralink 보고; 동료 심사 논문 없음 & 측정됨(회사 보고) \\
9 & 전극 100개 남짓의 배열을 쓴 인터페이스가 커서를 조작한다 & 마비 환자, 일차 운동 피질에 전극 $96$개: 손상 3년 뒤에도 팔을 움직이려는 의도가 스파이크를 바꾼다. 디코더가 “신경 커서”(메일, TV), 의수 손과 로봇 팔 조작을 가능하게 한다. & \cite{Hochberg2006} & 측정됨 \\
10 & “생각 속의 손”으로 글쓰기: 분당 약 $90$자, $8$~bit/s 미만 & 손이 마비된 사람, 손으로 쓰려는 시도, 순환 신경망 디코더: 실시간으로 분당 $90$자, 정확도 $94.1\,\%$, 자동 교정 뒤 $99\,\%$ 이상. 비트 추정은 본문의 산술. & \cite{Willett2021} & 측정됨 \\
11 & 말하려는 시도를 분당 약 $60$단어로 글로 바꾼다 & 또렷한 말을 잃은 사람, 피질의 배열: 분당 $62$단어; 어휘 $50$단어에서 단어 오류 $9.1\,\%$, 어휘 $125\,000$단어에서 $23.8\,\%$. & \cite{Willett2023} & 측정됨 \\
12 & 인터페이스는 용량의 작은 일부만 꺼낸다: 뉴런 하나는 초당 수십 비트, 100개는 수천 비트(명제~\ref{prop:probe}) & 상한~\eqref{eq:spikeentropy}은 이론; 8행, 10행과의 비교는 본문의 결론. 병목이 전선이 아니라 저차원성과 부호에 대한 무지라는 것은 직접 실험으로 확인되지 않았다. & \cite{Rieke1997} & 이론 \\
13 & 뇌와 디코더는 서로에게서 배운다 & 원숭이가 커서를 조작; 디코더가 닫힌 루프에서 맞춰지는 동안 뉴런도 다시 학습한다: 정확도가 오르고, 기술이 유지되며, 자기 팔의 움직임에 덜 방해받는다. 학습 속도를 떼어 놓으라는 조언은 공학 규칙이며, 본문에 별도 실험은 없다. & \cite{Orsborn2014coad} & 측정됨 \\
14 & 전극을 통한 전류는 주변 뉴런과 전선을 유형을 가리지 않고 흥분시킨다 & 생쥐와 고양이 피질, 2광자 칼슘 기록: 약한 전류조차 수십 마이크로미터 지름의 부피에서 축삭을 직접 흥분시켜, 최대 수 밀리미터 떨어진 드문드문한 뉴런을 흥분시킨다. 즉 흥분은 비선택적이지만 빽빽하지는 않다. & \cite{Histed2009stim} & 측정됨 \\
15 & 시각 피질 자극은 점을 켠다. 최대 $16$개 점을 동시에 켜면 몇몇 글자와 물체의 경계를 알아볼 수 있다 & 시각 장애인, 시각 피질에 전극 $96$개를 $6$개월 이식: 전극 하나의 문턱값 $66.8\pm36.5$~\textmu A; 무리를 동시에 자극하면(최대 $16$전극, 자극기의 한계) 문턱값이 낮아지고 구별 가능한 무늬가 생기며, 몇몇 글자와 물체의 경계를 알아본다. & \cite{Fernandez2021} & 측정됨 \\
16 & Neuralink의 두 번째 방향은 시각 피질에 신호를 넣는 장치다 & 개발 중이라는 회사 발표뿐. 작동에 관한 동료 심사 데이터는 찾지 못했다. & Neuralink 보고 & 가설 \\
17 & 브로드캐스트 신경전달물질의 농도는 조직 안의 화학 센서로 잴 수 있다 & 쥐 뇌 절편, 탄소 섬유 고속 순환 전압전류법: 세로토닌의 방출과 재흡수를 측정; 물질이 시냅스 밖으로 나간다. & \cite{Bunin1998} & 측정됨 \\
\bottomrule
\end{longtable}}

\section{\ref{sec:eetypes}장. 계측기로서의 뉴런}

개별 세포의 동작 모드는 전극을 통한 전류와 전압 기록으로 직접 측정되었다. 적분기의 수학과 부류 구분은 고전 이론이다. 뇌가 모드 재구성을 계산에 쓴다는 것은 여전히 가설이다.

{\footnotesize
\setlength{\tabcolsep}{3pt}\renewcommand{\arraystretch}{1.15}
\begin{longtable}{>{\raggedright}p{0.5cm}>{\raggedright}p{4.0cm}>{\raggedright}p{5.6cm}>{\raggedright}p{2.3cm}>{\raggedright\arraybackslash}p{1.7cm}}
\toprule
No. & 명제 & 실험과 측정한 것 & 출처 & 상태 \\
\midrule\endfirsthead
\toprule
No. & 명제 & 실험과 측정한 것 & 출처 & 상태 \\
\midrule\endhead
1 & 공진기: 전압 변화에 맞서는 느린 전류가 뉴런을 몇 Hz에서 수십 Hz에 최대가 있는 대역 통과 필터로 만든다(\ref{sec:resonator}절) & 절편의 \nt{GLUT} 뉴런에 주파수가 서서히 오르는 사인파 전류를 넣고 임피던스 $|Z(f)|$를 측정: $33^\circ$C에서 $2$--$5$~Hz에 최대($38^\circ$C에서 약 $7$~Hz). 공진은 느린 칼륨 전류와 양이온 전류가 만들고 지속성 나트륨 전류가 키운다. 리뷰는 같은 기법을 여러 세포 부류에 적용한다 & \cite{HuVervaekeStorm2002,HutcheonYarom2000} & 측정됨 \\
2 & 느린 전류는 인덕턴스처럼 행동하고, 막의 커패시턴스와 함께 공진 회로를 이룬다 & 전류에 대한 축삭의 문턱 아래 반응을 측정하고 선형화한 컨덕턴스 방정식으로 기술했다. 등가 회로에는 값이 전압에 따라 달라지는 “현상론적” 인덕턴스가 들어 있다 & \cite{Mauro1970} & 이론 \\
3 & 되먹임이 있는 무리의 시상수 $\tau_{\text{eff}}=\tau/(1-w)$~\eqref{eq:perfectint}; $\tau=0.1$~s, $w=0.995$에서 $20$~s & 본문에서 무리의 선형 방정식으로 유도. 눈 위치 유지의 메커니즘으로는 이론 연구에서 제안되었다 & 본문에서 유도; \cite{Seung1996} & 이론 \\
4 & 눈 위치 적분기는 개별 세포의 성질이 아니라 네트워크로 입력을 유지한다 & 물고기에서 눈 위치를 유지하는 핵의 세포 안 기록: 빠르게 눈을 돌린 뒤 뉴런의 발화율이 계단처럼 바뀌어 유지된다. 세포 하나에 짧은 전류를 넣으면 잠깐 이동할 뿐이고, 전압 계단은 문턱 아래에서도 유지된다. 이를 유지하는 것은 시냅스 입력이다 & \cite{Aksay2001} & 측정됨 \\
5 & 무누설 적분기에는 학습에 의한 끊임없는 조정이 필요하다. 어긋나면 잊거나 포화로 치닫는다 & 물고기를 눈 위치의 $\pm$에 비례하는 시각 되먹임으로 훈련: 유지가 불안정해지거나 누설되고, 뉴런의 시상수가 한 자릿수 넘게, $<1$~s까지 떨어졌다. 평범한 시각 되먹임은 안정성을 서서히 되돌렸다 & \cite{Major2004} & 측정됨 \\
6 & 쌍안정 뉴런: 짧은 입력이 오래가는 플래토를 켜고 억제가 끈다. 이 성질은 세로토닌이 켠다(\ref{sec:bistable}절) & 고양이에서 근육을 제어하는 \nt{ACET} 뉴런의 세포 안 기록: 짧은 시냅스 흥분이 뉴런을 오래가는 발화로 옮기고 짧은 억제가 재설정한다. 척수를 절단한 고양이에서는 세로토닌 전구체를 주입한 뒤 이 상태가 나타난다 & \cite{Hounsgaard1988} & 측정됨 \\
7 & 두 부류: 적분기(발화율이 0부터 커진다)와 공진기(문턱에서 곧바로 유한한 발화율) & 부류는 분기 이론에서 유도되었다. 실험: 피질 절편에서 \nt{GLUT} 뉴런은 얼마든지 낮은 발화율로 발화를 시작하고(1형), 빠른 \nt{GABA} 뉴런은 유한한 발화율로 도약한다(2형) & \cite{Izhikevich2007,Tateno2004} & 측정됨 \\
8 & 같은 뉴런이 적분기 또는 동시성 검출기가 된다. 억제가 합산 창을 짧게 한다 & 해마 절편에서 흥분성 입력이 문턱까지 더해지는 창은 $2$~ms보다 짧다. 흥분 직후에 오는 피드포워드 억제가 창을 짧게 한다. 억제가 약한 가지돌기에서는 창이 넓다 & \cite{Pouille2001} & 측정됨 \\
9 & 축삭으로 된 지연선(표~\ref{tab:eetypes}) & 원숭이올빼미에서 동시성 검출기로 가는 축삭의 지연을 측정: 축삭 길이가 다르면 지연이 다르고, 검출기는 소리 도달 시간차에 맞춰져 있다 & \cite{Carr1990} & 측정됨 \\
10 & 깨어 있을 때는 페이스메이커, 잘 때는 조용한 세포 & \nt{SERO} 뉴런은 낮에는 거의 규칙적으로 발화하고, 서파 수면에서 느려지며, 렘수면에서는 거의 조용하다(자세한 내용은 \ref{sec:sleep}장) & \cite{JacobsAzmitia1992,McGintyHarper1976} & 측정됨 \\
11 & 소자는 이온 채널과, 그 채널을 조정하는 브로드캐스트 채널이 정한다(명제~\ref{prop:eetypes}) & 작은 네트워크에서 입증: 조절 물질이 뉴런의 고유 성질과 시냅스 세기를 바꿔, 같은 회로가 서로 다른 리듬을 낸다 & \cite{Marder2012} & 측정됨 \\
12 & 뇌는 모드 재구성을 계산에 쓴다(명제~\ref{prop:eetypes}) & 과제를 푸는 데 세포 모드의 재구성이 필요했다는 직접 실험은 없다. 조절 물질이 회로의 출력을 바꾸는 실험만 있다 & \cite{Marder2012} & 가설 \\
\bottomrule
\end{longtable}}

\section{\ref{sec:dendrites}장. 가지돌기 개수}

부류의 구조와 입력 개수는 해부학과 전기 기록으로 측정되었다. 추정~\eqref{eq:dendriteload}은 커패시터의 간단한 물리이고, 입력 개수에 대한 계산적 설명은 이론과 모델에 기댄다.

{\footnotesize
\setlength{\tabcolsep}{3pt}\renewcommand{\arraystretch}{1.15}
\begin{longtable}{>{\raggedright}p{0.5cm}>{\raggedright}p{4.0cm}>{\raggedright}p{5.6cm}>{\raggedright}p{2.3cm}>{\raggedright\arraybackslash}p{1.7cm}}
\toprule
No. & 명제 & 실험과 측정한 것 & 출처 & 상태 \\
\midrule\endfirsthead
\toprule
No. & 명제 & 실험과 측정한 것 & 출처 & 상태 \\
\midrule\endhead
1 & 막의 단위 면적당 커패시턴스 $c_m\approx1$~\textmu F/cm$^2$ & 뉴런 세포체에서 막을 피펫으로 떼어 내고 전압 계단을 걸어, 용량성 전류의 감쇠로 전체 커패시턴스를 구한 뒤 면적으로 나눴다: 세 부류의 뉴런에서 $0.9$~\textmu F/cm$^2$ & \cite{Gentet2000} & 측정됨 \\
2 & 충전 시간 $t\approx c_mA\Delta V/I$~\eqref{eq:dendriteload}: 큰 뉴런에는 동시 입력 수십 개가 필요하고, 작은 뉴런에는 큰 시냅스 하나면 충분하다 & 커패시터 충전 법칙에 의한 추정. 수치($20$과 $300$~pF, $0.1$과 몇 nA)는 크기 수준 & 본문에서 유도 & 이론 \\
3 & 거대 시냅스 하나는 몇 나노암페어의 전류를 주어 작은 뉴런을 곧바로 문턱값까지 올린다 & 절편에서 말단과 수신 세포를 동시 기록: 시냅스 하나의 전류 $2$--$13$~nA, 입력 스파이크마다 문턱 이상의 반응, $36^\circ$C에서 지연 약 $0.4$~ms; 수신 세포는 \nt{GLYC}를 낸다 & \cite{Borst1995,BorstSoria2012} & 측정됨 \\
4 & 가지돌기 0개: 촉각과 통증의 감각 뉴런에서 스파이크는 세포체를 거치지 않고 센서에서 척수로 간다 & 구조(세포체 옆에서 둘로 갈라지는 축삭)는 해부학으로 확립되었다. 전도에 세포체의 흥분성이 필요 없다는 것은 이런 뉴런의 검증된 모델에서 보였다: 흥분성 세포체가 없어도 스파이크는 분기점을 똑같이 확실하게 통과한다 & \cite{AmirDevor2003} & 이론 \\
5 & 가지돌기 1개: 후각 뉴런은 수용체 한 종류를 지니고 입력 선로 하나를 전달한다 & 수용체 유전자 실험의 리뷰: 후각 뉴런 하나는 큰 유전자군 가운데 수용체 유전자 하나를 발현한다 & \cite{Mombaerts2004} & 측정됨 \\
6 & 가지돌기 2개: 소리 방향 검출기에서 왼쪽 귀와 오른쪽 귀의 입력은 서로 다른 가지돌기에 붙는다 & 리뷰에 모은 포유류의 해부학과 기록: 두 귀의 입력이 반대쪽을 향한 두 가지돌기에 나뉘어 있다 & \cite{Grothe2010} & 측정됨 \\
7 & 입력을 두 가지돌기에 나누면 AND가 더 엄격해진다 & 모델에서 입증: 한 가지돌기의 포화~\eqref{eq:shunt} 때문에 한쪽 입력만으로는 문턱값에 이르지 못하고, 동시성 검출이 개선된다. 입력 배치를 바꾼 직접 실험은 없다 & \cite{AgmonSnir1998} & 이론 \\
8 & 운동 학습 회로에 작은 \nt{GLUT} 뉴런 약 $7\cdot10^{10}$개, 각각 입력 $\sim4$개 & 등방성 분획법에 의한 세포 계수: 사람 뇌의 이 부분에 뉴런 약 $690$억 개. 살아 있는 쥐에서 이런 세포 하나를 기록: 촉각에 소수의 입력에서 오는 이산 전류의 버스트가 도착하고, 이들이 더해지면 세포가 발화한다 & \cite{Azevedo2009,Chadderton2004} & 측정됨 \\
9 & 입력 $n$개인 문턱 소자는 무작위 패턴을 최대 $P_{\max}\approx2n$개 분리한다~\eqref{eq:cover}; 운동 학습 회로의 출력 뉴런에서 이 경계는 $\sim2\cdot10^5$, 현실적인 추정은 $\sim4\cdot10^4$개의 대응 & 경계는 조합 기하학의 고전적 결과. 현실적인 추정은 양의 가중치만 쓰고 잡음 여유를 둔 용량 이론을, 이 뉴런의 측정된 입력 수에 적용한 것이다 & \cite{Cover1965,Brunel2004} & 이론 \\
10 & 입력이 적은 믹서는 입력을 확장하려고 필요하다. “네 개”는 최적에 가깝다 & 이론: 믹서 층이 주는 표현의 차원은 해부학적으로 관찰되는 입력 개수 근처에서 최대가 된다. 확장이라는 원래 가설은 고전 연구에 있다 & \cite{Marr1969,Albus1971,LitwinKumar2017} & 가설 \\
11 & 큰 나무: 입력 $10^4$--$10^5$개 & 전자 현미경 사진으로 시냅스 계수: 쥐의 운동 학습 회로 출력 \nt{GABA} 뉴런 하나에 시냅스 $\approx1.7\cdot10^5$개; 해마 \nt{GLUT} 뉴런에 흥분성 입력 약 $30\,000$개와 억제성 입력 $1\,700$개 & \cite{NapperHarvey1988,Megias2001} & 측정됨 \\
12 & 큰 나무의 가지는 각자 문턱값을 가진 별도의 하위 회로이며, 뉴런은 작은 다층 네트워크다 & 가는 가지돌기의 두 지점을 국소 자극: 같은 가지의 입력은 시그모이드로, 다른 가지의 입력은 선형으로 더해진다. 사람 피질 뉴런의 가지돌기에서 세포 하나가 선형 분리 불가능한 문제를 풀게 하는 등급형 칼슘 스파이크가 발견되었다. 모델: 뉴런 하나를 재현하려면 깊은 네트워크가 필요하다 & \cite{Polsky2004,Gidon2020,Beniaguev2021} & 측정됨 \\
13 & 출력 가지돌기: 망막 \nt{GABA} 뉴런의 가지 하나하나가 별도의 방향 검출기다 & 개별 가지의 2광자 칼슘 기록: 가지의 신호는 세포체에서 끝 쪽으로 가는 움직임에 선택적이지만 세포체의 전압은 그렇지 않다 & \cite{Euler2002} & 측정됨 \\
14 & 입력 개수는 과제를 따른다(명제~\ref{prop:dendrites}) & 부류의 구조는 측정되었다(3--13행). 입력 개수가 속도나 분류를 위해 선택되었다는 것은 개별 부류에 대해 이론과 모델로만 보였다 & \cite{LitwinKumar2017,AgmonSnir1998} & 가설 \\
\bottomrule
\end{longtable}}

\section{\ref{sec:sleep}장. 수면}

수면 모드, 재생, 시냅스 축소는 뉴런 기록, 미세투석, 전자 현미경으로 측정되었다. 수면 일정과 느린 시소는 이 책의 모델로 기술했고, 수면의 목적은 대부분 가설이다.

{\footnotesize
\setlength{\tabcolsep}{3pt}\renewcommand{\arraystretch}{1.15}
\begin{longtable}{>{\raggedright}p{0.5cm}>{\raggedright}p{4.0cm}>{\raggedright}p{5.6cm}>{\raggedright}p{2.3cm}>{\raggedright\arraybackslash}p{1.7cm}}
\toprule
No. & 명제 & 실험과 측정한 것 & 출처 & 상태 \\
\midrule\endfirsthead
\toprule
No. & 명제 & 실험과 측정한 것 & 출처 & 상태 \\
\midrule\endhead
1 & 잠자는 동안 피질 뉴런은 조용하지 않다. 바뀌는 것은 동작 모드다 & 쥐 피질 뉴런의 장시간 기록: 서파 수면에서 발화율은 0이 아니며, 활동 기간과 침묵 기간이 번갈아 온다. 오래 깨어 있은 뒤에는 모든 상태에서 발화율이 높고, 잔 뒤에는 낮다 & \cite{Vyazovskiy2009} & 측정됨 \\
2 & \nt{ACET}은 낮과 렘수면에서 높고 서파 수면에서 낮다(표~\ref{tab:sleepmodes}) & 자유롭게 움직이는 고양이의 피질 미세투석: 각성 시 아세틸콜린 방출이 서파 수면보다 $100$--$175\%$ 높고, 렘수면에서는 조용한 각성과 거의 같은 수준이다 & \cite{Marrosu1995} & 측정됨 \\
3 & \nt{NORA}와 \nt{SERO}는 서파 수면에서 잠잠해지고 렘수면에서 거의 조용하다 & 쥐의 개별 \nt{NORA} 뉴런과 고양이의 \nt{SERO} 뉴런을 수면 단계별로 기록: 발화율은 각성에서 최대, 서파 수면에서 낮고, 렘수면에서 거의 0이다 & \cite{AstonJonesBloom1981,McGintyHarper1976} & 측정됨 \\
4 & 렘수면에서 근육은 \nt{GABA}와 \nt{GLYC}를 통한 억제로 꺼진다 & 쥐에서 씹기 근육의 \nt{ACET} 뉴런 활동을 재고 차단제를 그 뉴런에 직접 주입: 두 유형의 \nt{GABA} 수용체와 \nt{GLYC} 수용체를 모두 차단해야만 마비가 풀린다 & \cite{BrooksPeever2012} & 측정됨 \\
5 & 수면 압력 $S$는 문턱값이 두 개인 커패시터다~\eqref{eq:sleeppressure}, $\tau_r=18.2$~h, $\tau_d=4.2$~h & 2과정 모델. 시상수는 EEG 서파 파워가 각성에 따라 커지고 수면 중에 줄어드는 방식에서, 문턱값의 모양은 자발적으로 깨는 시각에서 얻었다 & \cite{Borbely1982,Daan1984} & 이론 \\
6 & 기계는 스스로 $16.1$~h 각성과 $8.0$~h 수면에 이른다(그림~\ref{fig:twoprocess}) & 문턱값이 오르내리는 \eqref{eq:sleeppressure}에 따른 계산, 스크립트 \texttt{models/sleep\_models.py}: $16.1$~h 각성과 $7.9$--$8.0$~h 수면 & --- & 모델 \\
7 & $40$~h 동안 자지 않은 뒤 $S=0.90$이지만 수면은 $8.8$~h뿐이다. 빚은 길이가 아니라 깊이로 갚는다 & 모델: \texttt{models/sleep\_models.py}가 $S=0.90$과 $8.8$~h를 준다. 실험: $40.5$~h 동안 자지 않은 여덟 명에서 첫 회복 밤의 EEG 서파 파워가 평소보다 유의하게 높다 & \cite{Borbely1981} & 모델 \\
8 & 물리적으로 $S$는 아데노신이다 & 고양이 미세투석: 각성을 조절하는 영역 가운데 하나에서 세포 밖 아데노신이 깨어 있는 동안 오르고, 수면 박탈로 더 오르며, 자는 동안 떨어진다. $S$가 아데노신뿐이라는 것은 입증되지 않았다 & \cite{PorkkaHeiskanen1997,Dunwiddie2001} & 가설 \\
9 & 서파 수면에서 피질은 시소처럼 오간다: 수백 ms의 활동과 침묵, 1초에 한 번보다 드물게($<1$~Hz, 마취 상태에서 대개 $0.3$--$0.4$~Hz) & 마취한 고양이의 세포 안 기록: $0.8$--$1.5$~s의 탈분극과 긴 과분극, 리듬 $<1$~Hz, 대개 $0.3$--$0.4$~Hz. 같은 활동-침묵 교대가 자연 수면에서도 보인다(1행) & \cite{Steriade1993,Vyazovskiy2009} & 측정됨 \\
10 & 시소는 되먹임과 피로가 있는 무리다~\eqref{eq:updown}. $b=5$에서 주기 $1.1$~s(모델의 값), 활동은 시간의 약 $35\%$; $b=1$에서 고른 활동(그림~\ref{fig:updown}) & 계산, 스크립트 \texttt{models/sleep\_models.py}: 주기 $1.11$~s, 활동 시간 비율 $0.35$(정수 개 주기에 대한 평균); $b=1$에서 활동 비율 $1.0$. 모델의 주기는 측정값보다 짧다(9행) & --- & 모델 \\
11 & \nt{ACET}은 시소를 끄고 피질을 고른 작동으로 옮긴다 & 고양이에서 \nt{ACET} 뉴런 핵을 자극하자 피질 뉴런의 $79\%$에서 느린 리듬이 차단되고 고른 작동으로 바뀌었다. 주로 긴 과분극이 사라졌기 때문이다. 아세틸콜린은 피로를 만드는 느린 칼륨 전류를 억누른다 & \cite{SteriadeAmzica1993,McCormick1992} & 측정됨 \\
12 & $7$--$15$~Hz 리듬 폭발은 피질과 중계 핵 사이의 \nt{GLUT}--\nt{GABA} 루프가 만든다 & 페럿의 시각 중계 핵 절편에서, 폭발은 중계 세포가 스스로 흥분시키는 이웃 \nt{GABA} 핵의 억제에 대해 내는 반동 버스트로 생긴다 & \cite{vonKrosigk1993,SteriadeMcCormickSejnowski1993} & 측정됨 \\
13 & 하루의 재생은 빨리 감기로 진행된다: 해마에서 약 $20$배, 피질에서 $6$--$7$배 빠르다 & 서파 수면에서 해마 뉴런의 반복 순서가 시간적으로 압축되어 나타난다. 트랙을 달린 뒤 장소 뉴런의 순서가 같은 순서로, 약 $20$배 압축된 $\sim100$~ms의 짧은 폭발로 반복되었다. 피질에서는 $6$--$7$배 압축 & \cite{Nadasdy1999,LeeWilson2002,Euston2007} & 측정됨 \\
14 & 재생과 피질의 맞물림을 강화하면 기억이 좋아진다(인과 관계) & 쥐에서 학습 뒤 재생에 맞춘 자극으로 재생과 피질의 서파 및 리듬 폭발의 연결을 강화: 다음 날 기억이 높았고, 대조군은 우연 수준이었다 & \cite{Maingret2016} & 측정됨 \\
15 & 재생의 목적은 느린 기억의 섞어 배우기다 & 두 학습 시스템 이론과 인공 네트워크 계산: 순차 학습은 옛것을 망가뜨리고, 섞어 배우기는 그렇지 않다. 재생이 바로 이를 위해 필요하다는 뇌 직접 실험은 없다 & \cite{McClelland1995} & 가설 \\
16 & 잔 뒤 시냅스는 크기에 비례해 줄어들고, 큰 시냅스는 거의 바뀌지 않는다 & 생쥐 피질 시냅스 $6920$개의 3차원 전자 현미경: 잔 뒤 접촉 면적이 깨어 있은 뒤보다 $\sim18\%$ 작고, 감소는 크기에 비례하며, 크고 안정된 시냅스는 바뀌지 않는다 & \cite{deVivo2017} & 측정됨 \\
17 & 수면의 주된 과제는 가중치 재정규화다 & 시냅스 항상성 가설. 1행과 16행은 이와 부합하지만 이것이 주된 기능임을 증명하지는 않는다 & \cite{TononiCirelli2014} & 가설 \\
18 & 렘수면은 세계 모델의 벤치 시험이다 & 모드(표~\ref{tab:sleepmodes})는 측정되었다. 네트워크가 세계 모델을 돌린다는 것은 이론적 추측이며 실험은 없다 & \cite{Hobson2009} & 가설 \\
19 & 수면 중 뇌 세척: 열린 문제 & 한 실험: 잠자는 동안 세포 사이 공간이 $60\%$ 넓고 청소가 빠르다. 측정 방법이 다른 다른 실험: 잠자는 동안과 마취 상태에서 청소가 눈에 띄게 느리다. 결과가 서로 모순된다 & \cite{Xie2013,Miao2024} & 가설 \\
20 & 국소 수면: 깨어 있는 동물의 피질 일부가 잠깐 침묵으로 빠진다 & 오래 깨어 있은 쥐에서 피질 한 영역의 뉴런이 서파 수면에서처럼 잠깐 조용해지고, 국소 EEG에 서파가 나타난다. 그런 순간 쥐는 과제에서 더 자주 실수한다 & \cite{Vyazovskiy2011} & 측정됨 \\
\bottomrule
\end{longtable}}

\section{\ref{sec:livingmachine}장. 살아 있는 뉴런으로 만든 기계}

전극 배열 위 배양의 행동은 기록과 자극을 쓰는 직접 실험으로 측정되었다. 일제 발화의 메커니즘은 시냅스 고갈 모델과 미세 배양 실험에 근거하고, 접시 안의 도파민에 관한 주장은 단 몇 건의 실험에 근거하며, 그중 일부는 저자들 스스로도 시연에 불과하다고 본다.

{\footnotesize
\setlength{\tabcolsep}{3pt}\renewcommand{\arraystretch}{1.15}
\begin{longtable}{>{\raggedright}p{0.5cm}>{\raggedright}p{4.0cm}>{\raggedright}p{5.6cm}>{\raggedright}p{2.3cm}>{\raggedright\arraybackslash}p{1.7cm}}
\toprule
No. & 명제 & 실험과 측정한 것 & 출처 & 상태 \\
\midrule\endfirsthead
\toprule
No. & 명제 & 실험과 측정한 것 & 출처 & 상태 \\
\midrule\endhead
1 & 칩 위의 배양: 대부분 \nt{GLUT}, 표본에 따라 다른 더 적은 비율의 \nt{GABA}. 해마 배양에서는 약 $6\%$ & 전극 배열 위의 뉴런 수천 개 네트워크는 표준 실험이다(3행). 해마 배양의 \nt{GABA} 뉴런 비율은 GABA 염색으로 세었다. 뉴런의 약 $6\%$. 피질 배양에 대해서는 검증된 숫자를 제시하지 않는다 & \cite{Benson1994} & 측정됨 \\
2 & 오가노이드: 인간 줄기세포에서 키운 1밀리미터 정도의 3차원 신경 조직 & 줄기세포에서 여러 뇌 영역을 가진 3차원 오가노이드를 키웠다. 전구세포 층과 성숙한 뉴런이 있는 피질도 포함된다 & \cite{Lancaster2013} & 측정됨 \\
3 & 입력이 없으면 배양은 배양에 따라 1초에서 몇 분 간격으로 일제 발화한다 & 밀도 $3\,000$--$50\,000$ 뉴런의 피질 배양 $58$개를 5주 동안 기록했다(30분 기록 $963$건). 2주 뒤에는 거의 모든 배양에서 네트워크 전체의 일제 발화가 우세하다. 일제 발화 간격은 $1$--$300$~s이며 표본에 크게 의존한다 & \cite{Wagenaar2006} & 측정됨 \\
4 & 일제 발화는 신경전달물질 고갈로 끝나고, 간격은 $T_{\text{burst}}\approx\tau_{\text{rec}}\ln\frac{1}{1-x^*}$~\eqref{eq:burstinterval}. $2$--$4$~s는 $\tau_{\text{rec}}=0.8$~s를 쓴 모델의 추정값이고, 측정된 회복은 더 느리다 & 공식은 본문에서 유도했다. 모델: 고갈되는 시냅스를 가진 네트워크는 스스로 네트워크 전체의 일제 발화를 만든다. 뉴런 $4$--$30$개의 미세 배양 실험: 일제 발화의 길이는 이전 일제 발화 뒤의 시간에 의존하고, 신경전달물질 소포를 고갈시키면 일제 발화가 사라진다. 저자들의 결론은 일제 발화가 신경전달물질 저장량에 의해 제한된다는 것이다. 거기서 저장량 회복은 수십 초이며, 같은 공식으로 수십 초에서 몇 분의 간격이 나온다 & 이 장의 유도; \cite{Tsodyks2000,CohenSegal2011,Tsodyks1997} & 이론 \\
5 & “입을 다무는 교사”: 원하는 응답이 나오면 자극을 끄고, 응답이 굳어진다 & 자극 후 $50\pm10$~ms에 선택한 응답이 나타날 때까지 $0.3$--$1$~Hz로 배양을 자극하고, 그 뒤 $5$~분 동안 자극을 껐다. 몇 주기 뒤에는 원하는 응답이 곧바로 나왔다. 학습 곡선을 그렸다 & \cite{Shahaf2001} & 측정됨 \\
6 & 배양이 테니스를 친다. 세션 안의 유의미한 랠리 증가는 완전한 되먹임일 때만 & 자세한 내용은 \ref{sec:dishbrain}장과 그 부록. 자극 대신 침묵을 주어도 세션 끝의 배양은 되먹임이 없을 때보다 낫지만 효과가 더 작다. 세션 간 학습은 불안정하다 & \cite{Kagan2022} & 측정됨 \\
7 & 배양이 자기 입력을 예측하도록 학습한다는 것은 가설. “지각력”에 관한 과장된 표현은 반박되었다 & 자유 에너지 원리는 이론적 제안이며, 더 단순한 설명과 구별하는 직접 실험은 없다. 반론 논문에서 연구자 30명이 루프 안의 행동 변화는 지능도 감각도 보여 주지 않는다고 지적했다 & \cite{Friston2010,Balci2023} & 가설 \\
8 & 알맞게 고른 자극 패턴으로 배양이 가상의 몸을 목표로 움직인다 & 프로그램이 현재 결과에 따라 훈련 자극 패턴을 스스로 골랐다. 수십 분 만에 네트워크는 지정된 상태에 도달했고, $2$~시간 훈련 뒤 가소성은 $80$~분 동안 유지되었다. 같은 자극을 행동과 무관하게 주면 효과가 없었다 & \cite{Bakkum2008} & 측정됨 \\
9 & 저수지: 학습하는 것은 선형 판독 $y=W_{\text{out}}\mathbf r$~\eqref{eq:reservoir}뿐 & 저수지 컴퓨팅 이론: 감쇠하는 기억을 가진 임의의 비선형 시스템에 학습된 선형 출력을 더한다 & \cite{Maass2002,JaegerHaas2004} & 이론 \\
10 & 저수지로 쓴 오가노이드가 약 $78\%$의 정확도로 화자를 구분한다(저자들의 보고). 오차로 학습하는 것은 판독이고, 오가노이드의 연결은 교사 없이 바뀐다 & 여러 화자의 말소리를 전극 배열의 전류 패턴으로 오가노이드에 넣었다. 학습된 판독은 약 $78\%$의 정확도로 화자를 구분했다고 저자들이 보고한다. 저자들은 훈련 중 오가노이드의 기능적 연결성이 바뀌었다고도 보고하고, 이를 오가노이드 안의 교사 없는 학습이라 부른다 & \cite{Cai2023} & 측정됨 \\
11 & 뉴런 100만 개는 스파이크에 약 $10^{-4}$~W를 쓴다~\eqref{eq:dishpower} & 추정: 피질 에너지 예산~\eqref{eq:energy}에서 얻은 스파이크당 비용 $10^{-10}$~J에 스파이크 수를 곱했다. 배양 전력의 직접 측정은 없다 & 이 장의 유도; \cite{Attwell2001} & 이론 \\
12 & 섬광으로 주는 도파민: 케이지 도파민 $1$~mM, $240$회 반복, 유발 스파이크 수 증가 & 원격 오가노이드 플랫폼에서 광반응성 케이지 속 도파민($1$~mM)을 자극이 더 많은 스파이크를 일으킬 때 자외선($365$~nm, $0.8$~s)으로 방출했다. $240$단계(약 $5$~시간) 동안 스파이크 수는 전체적으로 늘었다. 저자들은 실험 하나로는 도파민과의 관련을 확정할 수 없다고 쓴다. 대조 실험은 없다 & \cite{Jordan2024} & 가설 \\
13 & 살아 있는 세포의 도파민이 유기 트랜지스터의 가중치를 오래, 가역적으로 바꾼다 & 도파민을 분비하는 세포를 유기 뉴로모픽 소자 위에서 키웠다. 전극에서의 도파민 산화가 채널 전도도를 바꾼다. 가중치의 장기 변화와 그 복귀를 보였다 & \cite{Keene2020} & 측정됨 \\
14 & 작동하는 \nt{DOPA} 뉴런을 가진 오가노이드가 존재한다. 그 도파민을 교사로 쓴 적은 없다 & 인간 줄기세포 오가노이드에서 전기적으로 활성인 성숙한 \nt{DOPA} 뉴런과 도파민 생산이 확인되었다. 이 도파민이 다른 네트워크를 가르치는 실험은 문헌에서 찾지 못했다 & \cite{Jo2016} & 측정됨 \\
15 & 막대의 균형을 잡는 오가노이드: 프로그램이 고른 학습 자극이 무작위보다 낫고, 굳히기는 없으며, \nt{GLUT} 시냅스를 통한다 & 생쥐 피질 오가노이드를 “카트 위의 막대” 과제와 닫힌 루프로 연결했다. 대부분의 오가노이드에서 강화 학습이 고른 신호가 무작위 신호나 무신호보다 좋은 결과를 냈다. $45$~분 휴식 뒤에는 향상이 남지 않았다. AMPA 수용체와 NMDA 수용체를 차단하면 향상이 사라졌다 & \cite{Robbins2026} & 측정됨 \\
16 & 제어 레지스터가 없으면 테스트 벤치의 네트워크는 무엇을 성공으로 볼지 스스로 정하지 못한다(명제~\ref{prop:livingmachine}) & 5--15행의 모든 실험에서 학습 신호는 실험자나 프로그램이 정한다. 배양이 스스로 성공 기준을 만든 실험은 없다. 이것은 부재에서 나온 결론이지 실험이 아니다 & --- & 가설 \\
\bottomrule
\end{longtable}}

\section{\ref{sec:dishbrain}장. DishBrain}

테스트 벤치의 구조와 게임 결과는 실험 논문 하나와 그 후속 연구에서 가져왔다. 잡음 공식과 좋은 설정 쪽으로의 표류는 본문에서 유도하고 이 책의 모델로 확인했으며, 접시 안의 학습 메커니즘은 밝혀지지 않았다.

{\footnotesize
\setlength{\tabcolsep}{3pt}\renewcommand{\arraystretch}{1.15}
\begin{longtable}{>{\raggedright}p{0.5cm}>{\raggedright}p{4.0cm}>{\raggedright}p{5.6cm}>{\raggedright}p{2.3cm}>{\raggedright\arraybackslash}p{1.7cm}}
\toprule
번호 & 주장 & 실험과 측정된 것 & 출처 & 상태 \\
\midrule\endfirsthead
\toprule
번호 & 주장 & 실험과 측정된 것 & 출처 & 상태 \\
\midrule\endhead
1 & 테스트 벤치: 칩당 생쥐 피질 세포 약 $800\,000$개 또는 인간 세포 약 $10^6$개. $8$~mm$^2$에 전극 $26\,000$개, 기록 최대 $1024$, 자극은 $8$개로 & 논문의 방법: MaxOne 칩, $8$~mm$^2$에 백금 전극 $26\,000$개, 초당 $20\,000$ 표본으로 $1024$채널 기록. 자극은 기술적으로 $32$개 전극까지 가능하나 독립적으로는 $8$개. 생쥐 배양은 약 $10$일째부터 사용 & \cite{Kagan2022} & 측정됨 \\
2 & $10$~ms 클록의 루프: 운동 영역의 스파이크가 라켓을 움직인다(그림~\ref{fig:dishloop}) & 스파이크를 $10$~ms 창으로 센다. 스파이크에서 응답 자극까지의 지연은 약 $5$~ms로, 장비의 버퍼가 정한다 & \cite{Kagan2022} & 측정됨 \\
3 & 입력: 장소 부호($8$개 전극 중 어느 것)와 초당 $4$--$40$회의 빈도 부호 & 본 실험에서 자극 빈도는 공이 먼 벽에 있을 때 $4$~Hz에서 라켓 앞 $40$~Hz까지 선형으로 오른다. 공 펄스의 진폭은 $75$~mV, 펄스는 직사각형 2상 & \cite{Kagan2022} & 측정됨 \\
4 & 출력 $\Delta p=c\,(n_\uparrow-n_\downarrow)$~\eqref{eq:dishreadout}, 창당 스파이크 수의 잡음 $1/\sqrt{RT}$ & 두 영역의 차이 규칙은 논문에 기술되어 있으나(활동에 따른 영역 가중치 포함) 정확한 공식은 없다. 잡음 추정은 푸아송 계수의 결과이며 본문에서 유도했다 & \cite{Kagan2022}; 본문의 유도 & 이론 \\
5 & 교사: 히트 뒤에는 예측 가능한 자극 $75$~mV, 초당 $100$회, $0.1$~s. 미스 뒤에는 예측 불가능한 자극 $150$~mV, 초당 $5$회, $4$~s 동안 무작위 장소와 시점에, 그 뒤 $4$~s 휴지 & 방법: 히트 뒤 $75$~mV, $100$~Hz, $100$~ms를 $8$개 전극 모두에 동시에. 미스 뒤 $150$~mV, $5$~Hz를 $4$~s 동안 무작위 전극에 무작위 시점으로, 그 뒤 $4$~s 자극 없음. 배양에 도파민은 없다 & \cite{Kagan2022} & 측정됨 \\
6 & 네트워크는 자기 입력이 예측 가능해지도록 행동한다 & 자유 에너지 원리는 저자들이 프로토콜을 이끌어 낸 이론적 제안이다. 실험은 이와 맞지만 더 단순한 메커니즘(7--8행)과 구별하지 못한다 & \cite{Friston2010,Kagan2022} & 가설 \\
7 & 실패가 네트워크를 흔들고 성공은 흔들지 않으면, 설정에 머무는 시간은 $\rho\propto1/P_{\text{miss}}$~\eqref{eq:dishdensity}(명\ref{prop:dishscramble}) & 대칭적 흔들기를 가진 마르코프 연쇄의 흐름 균형에서 나오며, 본문에서 유도했다. 배양에서 직접 확인한 것은 없다 & 본문의 유도 & 이론 \\
8 & `미스 때 뒤섞기' 모델이 학습한다: 히트 비율 $0.21\to0.38$. 모든 랠리 뒤 흔들기는 $\approx0.12$, 흔들기 없음은 $0.19$(그림~\ref{fig:dishmodel}) & 계산, 스크립트 \texttt{models/dishbrain\_model.py}, 모델 배양 $40$개에 랠리 $2000$번씩: $0.21\to0.38$; $0.13\to0.11$; $0.19\to0.19$ & --- & 모델 \\
9 & 랠리가 길어지는 것은 완전한 루프의 살아 있는 배양뿐이고, 대조군은 학습하지 않는다 & $20$~분 세션 $399$번: 생쥐 배양 $9$개(세션 $101$번), 인간 $11$개($138$), 배지만 있는 칩 $6$개($80$), 휴지 배양 $20$개($42$), 모의실험 $3$개($38$). 마지막 $15$~분의 평균 랠리 길이가 처음 $5$~분보다 긴 것은 살아 있는 배양뿐이다. 뉴런이 아닌 세포(HEK293T)와 배지만 있는 칩에서는 효과가 없다 & \cite{Kagan2022} & 측정됨 \\
10 & 세션 안의 유의미한 향상은 완전한 피드백일 때만. 자극 대신 침묵을 주면 세션 끝의 게임은 피드백 없음보다 낫지만 효과가 더 작다 & $3$일 동안 세션 $486$번: `자극'($n=27$), `침묵'($17$), `피드백 없음'($15$). 시간에 따른 향상은 `자극' 조건에서만 유의미하다. 세션 끝에 `침묵' 조건은 더 작은 효과로 `피드백 없음'을 앞섰다 & \cite{Kagan2022} & 측정됨 \\
11 & 효과는 작다. 네트워크가 오래 가는 학습을 하는지는 열린 질문 & 세션 안의 향상은 확실하지만, 저자들 스스로 말하듯 며칠에 걸친 세션 간 학습은 안정적으로 관찰되지 않았다 & \cite{Kagan2022} & 측정됨 \\
12 & 게임 중 네트워크는 임계 상태에 다가가고, 가까울수록 게임을 잘한다 & 같은 배양의 활동 사태 분석: 구조화된 입력은 네트워크를 임계 상태 쪽으로 옮기고, 게임 성적은 임계 상태와의 근접도와 상관된다. 그러나 행동의 결과에 대한 정보 없이 임계성만으로는 학습에 부족하다 & \cite{Habibollahi2023} & 측정됨 \\
13 & 배양의 `지각력'은 보이지 않았다 & 연구자 30명의 반론 논문: 닫힌 루프 안의 행동 변화는 지능도 감각도 증명하지 않는다. 그것을 보여 주는 실험은 없다 & \cite{Balci2023} & 가설 \\
14 & 제안하는 네 번째 대조 실험: 미스 뒤에 같은 길이와 같은 에너지의 예측 가능한 자극(\ref{sec:dishspec}절) & 이런 실험은 수행되지 않았다. 이 책의 제안이다 & --- & 가설 \\
\bottomrule
\end{longtable}}

